% Transcription — fonds Alexandre Grothendieck, shelfmark 1, batch 5 (pages 81-100).
% Established from the facsimile archives/batches/1/batch-05.pdf (a mirror of
% https://grothendieck.umontpellier.fr/1.pdf), per the skill
% transcribe-grothendieck.
%
% Pass: Fable 5.1 (claude-fable-5-1), 2026-09-10 — first pass, unchecked against
% the pages by a human.
%
% All twenty pages are transcribed. \page{} numbers the position in the
% folder, as batches 1 to 4 did: the archivists' pencil numbers run two
% ahead, so page 81 here is pencilled 83 and page 100 is pencilled 102.
%
% The batch is two paragraphs of the same treatise, in the fast blue-black
% hand of batches 3 and 4, each with its own formula numbering.
%
%   Pages 81-92 — « 4. Inégalités de convexité et classes remarquables
%     d'opérateurs », formulas (1) to (46). Th. 1 in two forms: 1° the
%     majorisation lemma attributed to H. Weyl — α decreasing, partial sums
%     of α bounded by those of β, φ convex symmetric, then φ(α) ≤ φ(β) under
%     a) φ increasing or b) equal total sums — and 2° its multiplicative
%     form for f(e^{t_1},…,e^{t_n}) convex, with the products ρ_1…ρ_m ≤
%     σ_1…σ_m; its proof through affine minorants, permutations (7)-(8) and
%     the Abel summation (9); Corollaire 1a on the reciprocals (10); a
%     Remarque on truncated sequences; the examples Σ f(ρ_i), the Schatten
%     norms N (12)-(14), and Σ ρ_i^p in (15)-(16). Th. 2 (« premier théorème
%     de convexité dans les espaces de Hilbert ») f(|λ_i|) ≤ f(ρ_i) in (17),
%     with Corollaires 1-5: (18)-(19) for Σ f, (20)-(22) for Σ|λ_i|^p ≤
%     Σρ_i^p = ‖u‖_p and the Fredholm and Hilbert-Schmidt cases, (23) for
%     f not increasing in finite dimension through (24)-(26) ρ_1…ρ_n =
%     |det u|, (27)-(28) for a Schatten norm. Th. 3 (« deuxième théorème »)
%     f(ρ_i(vu)) ≤ f(ρ_i(v)ρ_i(u)) in (29), with (30)-(33), the Hölder
%     inequality for operators (34)-(38) ‖vu‖_s ≤ ‖u‖_p‖v‖_q, a corollaire 3
%     in the margin |Tr vu| ≤ Σρ_i(u)ρ_i(v), and a pasted corollaire 4
%     (39)-(40). Th. 4 (« troisième théorème ») φ(ρ_i(u+v)) ≤
%     φ(ρ_i(u)+ρ_i(v)) in (41), the norm ‖u‖_N = N(ρ_i(u)) in (43), its
%     triangle inequality (44), and (45)-(46) ‖u+v‖_p ≤ ‖u‖_p + ‖v‖_p.
%   Page 87 is a small brownish sheet, half written, redrafting the end of
%     page 85 with its own (13) and (14).
%   Pages 93-100 — « 5. Classes remarquables d'opérateurs dans l'espace de
%     Hilbert », formulas (1) to (17): Définition 1 of a Schatten norm on
%     R^n and on R^(N), the examples ‖·‖_p in (2)-(3), the extension (4)-(5)
%     of N to all sequences, the spaces ℓ^N with (6) ℓ^1 ⊂ ℓ^N ⊂ ℓ^∞ and
%     the countable subadditivity (7), their completeness, the closure ℓ_N of
%     R^(N) with (8) ℓ^1 ⊂ ℓ_N ⊂ c_0, the polar norm N° in (9), the pairing
%     (10) and (ℓ_N)' = ℓ^{N°}, a Prop. on the reflexivity of ℓ_N, then the
%     operator classes L^N(E,F) and L_N(E,F) defined by (11), the norm N(u),
%     (12)-(14) in the margin, the multiplicativity (15)-(17), the duality
%     (L_N(E,F))' = L^{N°}(F,E) with its proof through v = U|v| and a
%     finite-rank projector, Corollaires 1-2 on |u| and on u = h + ik, and a
%     Remarque characterising the norms N(u) axiomatically on E'⊗F.
%
% Cancellations: pages 81, 82, 83, 84, 85, 86, 87, 88, 89, 90, 91, 92, 93,
% 94, 95, 96, 97, 98, 99 and 100 all carry lines or blocks struck through;
% each is transcribed as \struck{} with one \note{} for the block. The
% cancelled block at the head of page 82 and the block on page 98 carrying
% the duality statement are framed and barred; the latter is taken up again
% on page 99. Red pencil frames formulas (20), (21), (22), (36), (38), (44)
% and (46) of § 4 and (32), (33) of its margin, underlines words on pages
% 85, 91, 93, 94, 97 and 99, and leaves question marks in the margins of
% pages 91, 93, 97 and 99 and « !! » on page 97. Long sideways notes stand
% in the left margins of pages 81, 85, 86, 89, 91, 96, 98 and 99; those of
% 86, 89, 91 and 98 carry numbered formulas and are transcribed as
% \marginal{}; those of 81, 85, 96 and 99 are largely lost. « petits
% caractères » in the margins of pages 84 and 88 is a typesetting
% instruction. Two slips are pasted on pages 90 and 91 and transcribed in
% place.
%
% Notation: ρ_i(u) the singular values, λ_i(u) the eigenvalues, ‖·‖_p the
% norms Σρ_i^p of the n° 3, as in batch 4; N a Schatten norm, N° its polar,
% written N' on pages 97-98 and N° from page 96 and on page 99 — both are
% kept as written. R^(N), with N underlined, is the space of finitely
% supported sequences, set \mathbf{R}^{(\mathbf{N})}; c_0 the null sequences.
% His gothic S for the symmetric group is \mathfrak{S}_n. « Sup » and « lim »
% are set as he writes them. E'⊗F is his notation for the finite-rank
% operators, kept, with \bar{e}_i ⊗ f_i for the rank-one ones.
%
% The dating is the inventory's own for the shelfmark. Its group,
% « MATHÉMATIQUES / RECHERCHE / Travaux », carries no date.
\documentclass[11pt,a4paper]{article}
% Le préambule commun à toutes les transcriptions du fonds Grothendieck.
%
% Il tient en une idée : **l'incertitude se compose, elle ne se résout pas.**
% Une transcription qui lisse un mot douteux a détruit la seule chose pour
% laquelle on la faisait — un lecteur doit pouvoir savoir, à chaque endroit,
% ce qui était sur le feuillet et ce qui a été ajouté. Les six macros qui
% suivent sont donc l'appareil critique, et non de la décoration.
%
% Les mêmes commandes sont reconnues par `scripts/render.mjs`, qui produit la
% vue de lecture du site. Une macro ajoutée ici doit l'être aussi là-bas,
% faute de quoi le rendu échouera bruyamment — ce qui est voulu : un rendu
% silencieusement faux est pire qu'un rendu absent.

\ProvidesPackage{grothendieck}[2026/08/08 Transcriptions du fonds Grothendieck]

\usepackage{amsmath,amssymb,amsthm}
\usepackage{mathrsfs}
\usepackage{stmaryrd}
% Les diagrammes commutatifs sont partout dans ce fonds : c'est la langue
% même de la géométrie algébrique de Grothendieck, et une transcription qui
% les rendrait en prose ne transcrirait rien.
\usepackage{tikz-cd}
% Français par défaut — c'est la langue des feuillets — mais l'anglais est
% chargé aussi : la traduction et le résumé sont des documents anglais, et la
% typographie française (espaces avant les deux-points) y serait une faute.
% Ces documents basculent par \selectlanguage{english} en tête de corps.
\usepackage[english,french]{babel}
\usepackage{fontspec}
\usepackage{geometry}
\geometry{a4paper,margin=25mm,marginparwidth=28mm}
\usepackage{marginnote}
\usepackage{xcolor}
% ulem, not soul. `soul`'s \st and \ul are built on a character-by-character
% scan that XeLaTeX's font machinery breaks: the document fails with an
% undefined control sequence rather than degrading. ulem does the same two jobs
% and is Unicode-engine safe. `normalem` stops it hijacking \emph, which the
% transcriptions use for Grothendieck's own emphasis.
\usepackage[normalem]{ulem}
\usepackage{hyperref}
\hypersetup{colorlinks=true,linkcolor=black,urlcolor=black,pdfborder={0 0 0}}

% --- Métadonnées du lot -----------------------------------------------------
% Reprises telles quelles par le script de rendu, qui les affiche en tête de
% la vue de lecture. Elles fixent ce que le fichier prétend couvrir : sans
% elles, une transcription flotte sans point d'attache dans un fonds de seize
% mille feuillets.

\newcommand{\folder}[1]{\gdef\grothFolder{#1}}
\newcommand{\batch}[1]{\gdef\grothBatch{#1}}
\newcommand{\leaves}[2]{\gdef\grothFirst{#1}\gdef\grothLast{#2}}
\newcommand{\foldertitle}[1]{\gdef\grothFoldertitle{#1}}
\gdef\grothFolder{?}\gdef\grothBatch{?}\gdef\grothFirst{?}\gdef\grothLast{?}\gdef\grothFoldertitle{}

% --- Le résumé --------------------------------------------------------------
%
% Il ouvre la lecture modernisée, sous le titre « Résumé », et il est composé
% en petit. Ce n'est pas un ornement : c'est le seul passage du document qui
% ne soit pas des mathématiques, et un lecteur doit pouvoir le reconnaître
% avant de l'avoir lu — puis le sauter s'il connaît déjà le sujet, ou s'y
% arrêter s'il ne le connaît pas. Le filet qui le suit marque la reprise du
% corps.

\newenvironment{resume}
  {\small\setlength{\parskip}{0.45em}}
  {\par\medskip\hrule\medskip}

% --- Mots-clés ---------------------------------------------------------------
%
% En anglais, en clôture du résumé de la lecture modernisée : le vocabulaire
% moderne sous lequel chercher ce que les feuillets construisent. Le manifeste
% du site (`npm run manifest`) les extrait pour étiqueter la cote entière —
% cette ligne est la source unique des tags, il n'y a pas de fichier à côté.
\newcommand{\keywords}[1]{\par\smallskip\noindent
  {\footnotesize\textsc{Keywords} --- \textit{#1}}\par}

% --- L'appareil critique ----------------------------------------------------

% Le numéro de feuillet, en marge. C'est l'ancre qui permet de régler un
% désaccord sur une formule en regardant *un* feuillet plutôt qu'une chemise
% de six cents. Le site s'en sert aussi pour tourner les pages du fac-similé
% quand on fait défiler la transcription.
\newcommand{\leaf}[1]{%
  \marginnote{\footnotesize\color{gray!70}\textsf{#1}}%
  \label{leaf:#1}\ignorespaces}

% Illisible : on ne devine pas. Un mot inventé qui se lit comme les autres est
% le pire résultat possible.
\newcommand{\ill}{\textcolor{red!60!black}{[\,\ldots\,]}}

% Lecture douteuse : le mot est donné, et signalé comme incertain.
\newcommand{\uncertain}[1]{\uline{#1}}

% Ajout de l'éditeur — une lettre manquante, un mot sous-entendu.
\newcommand{\add}[1]{\textcolor{blue!50!black}{[#1]}}

% Biffé par Grothendieck lui-même. Ce qu'il a rayé fait partie du document :
% une rature dit souvent où la pensée a changé de direction.
\newcommand{\struck}[1]{\sout{#1}}

% Note du transcripteur, sur l'état matériel du feuillet.
\newcommand{\note}[1]{\par\smallskip{\small\color{gray!60!black}\textit{[#1]}}\par\smallskip}

% Note marginale de Grothendieck, distincte du corps du texte.
\newcommand{\marginal}[1]{\par\smallskip\noindent\hspace{1em}%
  {\small\color{gray!60!black}#1}\par\smallskip}

% --- Titre ------------------------------------------------------------------

\AtBeginDocument{%
  \begin{center}
    {\large\bfseries Fonds Alexandre Grothendieck --- cote n\textsuperscript{o}~\grothFolder}\\[2pt]
    {\itshape\grothFoldertitle}\\[4pt]
    {\small feuillets \grothFirst--\grothLast\ (lot \grothBatch)}\\[2pt]
    {\footnotesize\color{gray!60!black}Transcription établie d'après le fac-similé numérisé
     par l'Université de Montpellier.\\
     Les lectures incertaines sont soulignées, les illisibles notées [\,\ldots\,],
     les ajouts entre crochets.}
  \end{center}
  \vspace{1em}}

\endinput


\folder{1}
\batch{5}
\pages{81}{100}
\dating{1953}
\watermark{Édition de démonstration}
\foldertitle{Analyse fonctionnelle : notes manuscrites (s.d.), lettre (1953)}

\begin{document}

\section{Inégalités de convexité et classes remarquables d'opérateurs}

\note{Le titre est de lui, en tête de la page 81, souligné : « 4. Inégalités de convexité et classes remarquables d'opérateurs, compléments ». Le dernier mot est très abrégé et la lecture en est douteuse. Les formules sont numérotées à partir de (1) et courent jusqu'à (48) sur les pages 81 à 92.}

\page{81}
4. Inégalités de convexité et classes remarquables d'opérateurs, \uncertain{compléments}

\emph{Th. 1} \note{le numéro est cerclé, ainsi que le « 1° » qui suit ; deux lignes sont biffées sous le titre} 1° \struck{La méthode de \ill{}} (H. Weyl). \emph{Lemme fondamental.} Considérons deux suites finies $(\alpha_i)_{i \leq n}$ et $(\beta_i)_{i \leq n}$ de nbs réels, la \uncertain{deuxième} \uncertain{décroissante} \note{deux mots en interligne, incertains} et \note{le « (2) » qui suit est cerclé}
\[
  \text{(1)}\qquad \alpha_1 \geq \alpha_2 \geq \cdots \geq \alpha_n
\]
\[
  \text{(2)}\qquad \alpha_1 + \cdots + \alpha_m \leq \beta_1 + \cdots + \beta_m \qquad \text{pour } m = 1, \ldots, n,
\]
\note{(1) est encadrée, (2) encadrée d'un trait double} soit $\varphi(t_1, \ldots, t_n)$ fonction numérique convexe et \emph{symétrique} sur $\mathbf{R}^n$. On suppose vérifiée l'une des deux conditions suivantes : a) $\varphi$ est croissante en chaque argument $t_i$ \struck{\ill{} croissante \ill{} dans $\mathbf{R}^n$} ; b) $\sum_{i=1}^n \alpha_i = \sum_{i=1}^n \beta_i$. Sous ces conditions, on a
\[
  \text{(3)}\qquad \varphi(\alpha_1, \ldots, \alpha_n) \leq \varphi(\beta_1, \ldots, \beta_n) .
\]
\struck{\ill{} le lemme suivant \ill{} \ill{} \ill{} immédiatement} \note{trois lignes lourdement biffées}

2° \struck{Lemme} \note{le « 2° » est cerclé} Soit $f(\rho_1, \ldots, \rho_n)$ \note{en interligne au-dessus de « fonction »} fonction \struck{numérique} \note{en interligne : « numérique »} de $n$ arguments $\rho_i > 0$, telle que $f(e^{t_1}, \ldots, e^{t_n})$ soit une fonction \emph{convexe} de $(t_1, \ldots, t_n) \in \mathbf{R}^n$. Soient $(\rho_1, \ldots, \rho_n)$ et $(\sigma_1, \ldots, \sigma_n)$ deux suites finies \note{en interligne : « finies »}, telles que \note{« (4) » et « (5) » cerclés}
\[
  \text{(4)}\qquad \rho_1 \geq \rho_2 \geq \cdots \geq \rho_n
\]
\[
  \text{(5)}\qquad \rho_1 \cdots \rho_m \leq \sigma_1 \cdots \sigma_m \qquad \text{pour tout } m = 1, \ldots, n .
\]
\note{les deux formules encadrées} On suppose vérifiée l'une des deux conditions suivantes : a) $f$ est fonction croissante de chaque argument $\rho_i$ \note{en interligne, biffé : « \struck{chaque argument $\rho_i$} » ; la ligne « fonction croissante de \ill{} » est reprise deux fois} ; b) $\rho_1 \cdots \rho_n = \sigma_1 \cdots \sigma_n \neq 0$. Sous ces conditions on a
\[
  \text{(6)}\qquad f(\rho_1, \ldots, \rho_n) \leq f(\sigma_1, \ldots, \sigma_n) .
\]
\note{le « (6) » est précédé d'un signe cerclé, peut-être un « (4) » corrigé}
En effet, si les $\rho_i, \sigma_i$ sont $> 0$, cet énoncé est manifestement équivalent à l'énoncé du \struck{th.} \uncertain{lemme} \note{« th. » en interligne au-dessus du mot biffé} \ill{} \struck{\ill{} (2°) \ill{}} sans le cas \ill{}, en posant $\alpha_i = \log \rho_i$, $\beta_i = \log \sigma_i$. \struck{\ill{} \ill{}, puisque} \struck{Dans le cas a),} \struck{l'ensemble des systèmes} \note{deux lignes biffées, dont les mots « dans le cas a) » en interligne} \ill{} \uncertain{systèmes} $(\rho_i), (\sigma_i)$ tels que $\rho_i \geq 0$, \struck{$\rho_i > 0$}, $\rho_1 \geq \rho_2 \geq \cdots \geq \rho_n \geq 0$, \struck{\ill{}} $\rho_1 \cdots \rho_m \leq \sigma_1 \cdots \sigma_m$ pour $m = 1, \ldots, n$, \uncertain{est} \uncertain{limite} \struck{\ill{}} \struck{$\rho_1 \cdots \rho_n = \sigma_1 \cdots \sigma_n$) \ill{} l'adhérence \ill{}} \note{deux lignes biffées} dans $\mathbf{R}^{2n}$ \struck{des \ill{}} \uncertain{du} \uncertain{système} analogue \struck{\ill{} de \ill{} systèmes \ill{}} $(\rho_i'), (\sigma_i')$ avec $\rho_i' > 0$, $\sigma_i' > 0$ (d'où \uncertain{le} \uncertain{corollaire} \note{« (2°) » en interligne} dans le cas a) pour $f$ \ill{} \uncertain{continue} \struck{\ill{}} $\sigma_i > 0$ pour tout $i$ ; \ill{} \ill{}). En effet, si \struck{\ill{}} \uncertain{alors} \ill{}, \ill{} les $\rho_i$ qui seraient nuls \struck{\ill{}} \uncertain{et} \uncertain{il} \uncertain{suffit} \uncertain{d'approcher} \uncertain{les} \uncertain{systèmes} \ill{} \struck{\ill{} aux conditions imposées, et \ill{}} \struck{\ill{} systèmes $(\rho_i') = (\rho_i')$, avec \ill{} $\rho_i' = \sigma_i' = \frac{1}{n}$ \ill{}} \struck{\ill{} dans ce cas ;} \struck{Si plus généralement, donc d'ici la conclusion \uncertain{cherchée}} \struck{\ill{} d'abord \ill{}} \note{les six dernières lignes de la page sont biffées d'un trait chacune ; un trait ondulé au crayon rouge court le long de la marge gauche de cette moitié de page}
\marginal{\ill{}} \note{deux notes dans la marge gauche, écrites de côté : l'une au crayon, cerclée, en haut, où l'on lit « $\alpha_i$ », « (1) », « \uncertain{Xavier} » ; l'autre à l'encre, en regard du 2°, qui commence par un astérisque, « \uncertain{Il} \uncertain{est} \ill{} », et où l'on lit « \uncertain{soit} \uncertain{d'abord} », « \uncertain{corollaire} » ; l'une et l'autre restent illisibles même redressées}

\page{82}
\struck{\uncertain{Vérifions} $(\rho_i)$, $(\sigma_i)$ par les \uncertain{systèmes} $(\rho_i' = \rho_i)$, $(\sigma_i + \varepsilon)$ \ill{} $(\rho_i'), (\sigma_i')$ \ill{} $k = $ \ill{} \ill{} satisfont \ill{} aux \uncertain{conditions} \ill{} \ill{} $\sigma_{n-1}' \neq 0$ (\ill{} pour \ill{}} \note{quatre lignes en tête de page, entourées d'un cadre et biffées}

Il suffit donc, \uncertain{dans} \uncertain{le} \uncertain{cas} \ill{}, \note{en interligne : « systèmes admissibles »} d'approcher \note{« d'approcher » en interligne, au-dessus d'un mot biffé} les systèmes $(\rho_i), (\sigma_i)$ \uncertain{par} \uncertain{des} \note{en interligne : « $\sigma_{n-1} = 0$, \ill{} », biffé} $(\rho_i'), (\sigma_i')$ avec $\sigma_n > 0$, et on prendra \struck{\ill{} \ill{}}

\struck{\ill{} \uncertain{des} $(\rho_i'), (\sigma_i')$ avec $\sigma_{n-1}' > 0$, \ill{} \ill{} \uncertain{soit} \ill{} \uncertain{pour} \ill{} $(\rho_i') = (\rho_i)$, $\sigma_i' = \sigma_i + \varepsilon$ \uncertain{pour} \uncertain{tout} $i \leq n-1$, $\sigma_n' = \sigma_n$ (\ill{} \ill{} \ill{} \uncertain{admissible}, \ill{} \ill{} $\prod \rho_i' = \prod \sigma_i' = 0$) \ill{} \uncertain{pour} \ill{} \ill{} $\sigma_{n-1} > 0$, $\sigma_n = 0$, \struck{\ill{}} \ill{} \ill{} \uncertain{pour} $\sigma_i' = \sigma_i$ pour $i \leq n-1$, $\sigma_n' = \sigma_n + \varepsilon$ \ill{} \ill{}. \uncertain{Dans} \uncertain{le} \uncertain{cas} \ill{}, \uncertain{si} \ill{} \ill{} \uncertain{les} \uncertain{systèmes} $(\rho_i'), (\sigma_i')$ \uncertain{sont} \uncertain{admissibles}, \uncertain{i.e.} \ill{} \ill{} \ill{} $\prod \rho_i = \prod \sigma_i$ \ill{} \ill{} \ill{} \uncertain{pas} $\prod \rho_i' = \prod \sigma_i'$. \uncertain{Mais} \ill{} \ill{} \ill{} \uncertain{les} $\rho_i$ \uncertain{nuls}, \ill{} \uncertain{obtient} \ill{} \uncertain{que} \struck{\ill{}} $\rho_i' \cdots \rho_j' \leq \sigma_i' \cdots \sigma_j'$ \ill{} \ill{} \uncertain{augmentation} \ill{} \uncertain{les} \ill{} $\rho_1' \cdots \rho_n' = \sigma_1' \cdots \sigma_n'$. \uncertain{Prenons} $\rho_i' = \rho_i$, $\sigma_i' = \sigma_i + \varepsilon$.} \note{tout ce bloc, qui occupe la moitié supérieure de la page, est entouré d'un cadre et barré de quatre longs traits obliques ; il reprend, pour le cas b), l'argument d'approximation de la page précédente}

Prouvons maintenant le lemme. Il est bien connu \uncertain{qu'une} \uncertain{fonction} \uncertain{convexe} \ill{} $\varphi$ sur $\mathbf{R}^n$ est \uncertain{dans} \uncertain{l'enveloppe} \note{« fonctions » en interligne, « \struck{\ill{}} » biffé} \uncertain{supérieure} $\varphi = \operatorname{Sup}_k L_k$ \note{le « $L_k$ » est corrigé, surchargé} où les $L_k$ sont des fonctions \emph{\uncertain{linéaires}} \uncertain{affines} \note{mot en interligne, souligné, incertain} dans $\mathbf{R}^n$ (ça résulte du fait que le \uncertain{sous-ensemble} de $\mathbf{R}^n \times \mathbf{R}$ \uncertain{au-dessus} \uncertain{du} \uncertain{graphe} de $\varphi$ est \uncertain{convexe} \uncertain{par} \uncertain{définition} \uncertain{des} \uncertain{fonctions} \uncertain{convexes}, et \uncertain{fermé} \uncertain{si} $\varphi$ est \uncertain{continue}, --- donc \uncertain{intersection} de \uncertain{demi-espaces} \uncertain{fermés}, dont on constate \uncertain{aussitôt} \uncertain{qu'ils} \struck{\ill{}} peuvent \note{« peuvent » en interligne} \uncertain{se} \uncertain{mettre} \uncertain{sous} \uncertain{la} \uncertain{forme} \note{trait ondulé au crayon rouge sous cette ligne} \ill{} \uncertain{d'équations} \ill{} $L_k(t) \geq 0$). \uncertain{De} \uncertain{plus}, si $\varphi$ est \struck{\ill{}} \uncertain{croissante} \note{« croissante » en interligne}, \struck{\ill{} \ill{}} \uncertain{les} $L_k$ \uncertain{le} \uncertain{sont}, \struck{\ill{}} \uncertain{i.e.} \ill{} \uncertain{pour} $L_k(t) = \sum_{i=1}^n a_i t_i + b_k$, \struck{\ill{}} $a_i \geq 0$ \note{la fin de la ligne est biffée : « \struck{\ill{} le voir} »}. En effet, si on avait $a_i < 0$, alors \struck{la fonction \ill{}}

\page{83}
fixant les arguments $t_j$ pour $j \neq i$, $L_k(t) = L_k(t_1, \ldots, t_n)$ \struck{serait \ill{} fonction \uncertain{linéaire} \uncertain{strictement} \uncertain{décroissante} de $s = t_i$, \ill{} \ill{} \uncertain{fonction} \uncertain{continue} \uncertain{croissante} $\varphi(t_1, \ldots, t_n)$, \ill{} \uncertain{indifférente} \ill{} \uncertain{impossible}, \ill{} \ill{} $\lambda(s)$ \ill{} \uncertain{si} $s \to -\infty$, \ill{} \uncertain{pour} \ill{} \ill{} \uncertain{pour} $s \to -\infty$.} \note{bloc de cinq lignes entouré d'un cadre et barré ; les mots « $s = t_i$ », « $\lambda(s)$ », « $s \to -\infty$ » sont sûrs, la prose ne l'est pas} \uncertain{on} \uncertain{aurait} $L(t) \to \infty$ \uncertain{si} $t_i \to -\infty$, \uncertain{tandis} \uncertain{que} $\varphi(t)$ \struck{\uncertain{resterait}} \note{en interligne} \uncertain{borné} \uncertain{pour} $t_i \to -\infty$, en \uncertain{contradiction} \uncertain{avec} $L(t) \leq \varphi(t)$ \uncertain{pour} \uncertain{tout} $t$. \struck{Comme \ill{}} \note{en interligne : « \uncertain{Si} \uncertain{maintenant} »} $\varphi$ \uncertain{est} \uncertain{aussi} \uncertain{symétrique}, \struck{$L_k(t_{\sigma 1}, \ldots, t_{\sigma n}) \leq \varphi(t_{\sigma 1}, \ldots, t_{\sigma n})$} \note{en interligne au-dessus : « ($L_k$ \uncertain{fonction} \ill{}) » et « $\sigma L \leq \varphi$ »} \uncertain{pour} \uncertain{tout} $\sigma \in \mathfrak{S}_n$ (\uncertain{pour} \ill{} \uncertain{des} \uncertain{permutations}, \uncertain{on} \uncertain{pose} $\sigma L(t_1, \ldots, t_n) = L(t_{\sigma^{-1} 1}, \ldots, t_{\sigma^{-1} n})$) \note{le $\sigma$ des indices porte un accent ou un exposant, lu $\sigma^{-1}$ ; incertain}, \uncertain{donc} \struck{\ill{} \ill{} \ill{}} \note{ligne biffée sous la parenthèse}
\[
  \text{(7)}\qquad \struck{\varphi(t)}\; \varphi = \operatorname{Sup}_k \operatorname{Sup}_{\sigma \in \mathfrak{S}_n} \sigma . L_k \struck{(t_1, \ldots, t_n)}
\]
\note{le second membre de (7) est biffé après « $\sigma . L_k$ »}

Il s'ensuit qu'il \ill{} \uncertain{pour} \uncertain{prouver} (3), \uncertain{il} \uncertain{suffit} \uncertain{de} \uncertain{prouver} \uncertain{que} \struck{$L = L_k$, \ill{}} \note{en interligne} $\sigma L_k(\alpha_1, \ldots, \alpha_n) \leq \varphi(\beta_1, \ldots, \beta_n)$ \note{ligne biffée puis reprise ; le « $(\alpha_1, \ldots, \alpha_n)$ » est ajouté en interligne} pour \struck{\ill{}} \uncertain{tout} $k$, et \uncertain{tout} $\sigma \in \mathfrak{S}_n$. \struck{\ill{}} \uncertain{Désignant} \note{« Désignant » en interligne} \uncertain{alors} $\sigma L_k$ \uncertain{par} $L$, \uncertain{on} \uncertain{est} \uncertain{ramené} \uncertain{à} \uncertain{prouver}
\[
  \text{(8)}\qquad \operatorname{Sup}_{\sigma \in \mathfrak{S}_n} (\sigma . L)(\alpha_1, \ldots, \alpha_n) \leq \operatorname{Sup}_{\sigma \in \mathfrak{S}_n} (\sigma L)(\beta_{\sigma 1}, \ldots, \beta_{\sigma n})
\]
\note{le second membre est écrit « $\operatorname{Sup}_{\sigma \in \mathfrak{S}_n} (\sigma L)(\beta_{\sigma 1}, \ldots, \beta_{\sigma n})$ », le $\sigma L$ surchargé d'un $L$ ; les indices $\sigma 1, \ldots, \sigma n$ sont surchargés et incertains}
(car le deuxième membre \uncertain{est} \ill{} $\leq \varphi(\beta_1, \ldots, \beta_n)$). Dans cette \note{« Dans cette » en interligne, au-dessus de « \struck{Dans} \struck{\ill{}} », souligné au crayon rouge} \uncertain{inégalité}, $L(t_1, \ldots, t_n) = \sum_{i=1}^n a_i t_i + b$ \note{ligne soulignée} \struck{\ill{}} \uncertain{avec} \uncertain{des} \struck{\ill{}} \uncertain{coefficients} \uncertain{a\_i} \ill{} \uncertain{quelconques} \uncertain{dans} \uncertain{le} \uncertain{cas} b), et \uncertain{des} $a_i \geq 0$ \uncertain{pour} \uncertain{tout} $i$ \uncertain{dans} \uncertain{le} \uncertain{cas} \uncertain{a}). \struck{\ill{} \ill{} \ill{}} \note{deux lignes biffées, avec en interligne « On peut d'ailleurs \ill{} »} \struck{\ill{} \uncertain{il} \uncertain{suffit} \ill{}} \uncertain{supposer} $b = 0$. \struck{D'où \ill{}} \struck{$\sum a_i b_i$ \ill{}} \struck{$L(\alpha_1, \ldots, \alpha_n) \leq L(\ldots)$ \ill{} \uncertain{lorsque}} \note{trois lignes biffées}
\[
  \sum a_i \alpha_i \leq \sum a_i \alpha_{\sigma i} \qquad \struck{\ill{}} \ill{}
\]
\note{la formule et la ligne qui la suit sont barrées d'un grand trait ondulé ; on lit dans le trait « \uncertain{quitte} \uncertain{à} $\sigma$ \ill{} \uncertain{une} \uncertain{permutation} \ill{} »} \uncertain{telle} \uncertain{que} \uncertain{la} \uncertain{suite} \uncertain{des} $a_{\sigma i}$ \uncertain{soit} \uncertain{décroissante} \note{en interligne, plus petit : « \uncertain{visible} \uncertain{aussitôt} \uncertain{puisque} \uncertain{que} $a_1 \geq a_2 \geq \cdots$ \ill{} »} \ill{} \uncertain{ce} \uncertain{lemme} \uncertain{fondamental}, \uncertain{sans} \ill{} \ill{} D'\uncertain{où}, \ill{} $= \sum a_i \alpha_{\sigma^{-1} i}$ \note{en interligne au-dessus de la formule suivante}
\[
  \sum a_{\sigma i}\, \alpha_i \quad (\ill{} \uncertain{maximum} \ill{}), \qquad \text{pour } \sigma \in \mathfrak{S}_n \text{ variable},
\]
\note{« variable » incertain}

\page{84}
\uncertain{prend} \uncertain{le} \uncertain{max} $(a_{\sigma i})$ \uncertain{est} \uncertain{décroissante}, \struck{\ill{} \uncertain{on} \uncertain{peut} \uncertain{donc} \uncertain{se} \uncertain{ramener} \uncertain{au} \uncertain{besoin} \ill{}} \struck{\ill{} \ill{} $(\sigma)$, \ill{}} \struck{\ill{} pour \ill{} le \ill{} $a_1 \geq a_2 \geq \cdots \geq a_n$} \note{trois lignes biffées} \uncertain{ramené} \uncertain{à} \uncertain{prouver} \uncertain{que} \uncertain{si} $a_1 \geq a_2 \geq \cdots \geq a_n$, \uncertain{on} \uncertain{a} \struck{$\sum a_i \alpha_i$} \note{en interligne, biffé}
\[
  \sum a_i \alpha_i \leq \sum a_i \beta_i ,
\]
\note{le second membre est surchargé : un « $\sum a_i \alpha_i$ » corrigé en « $\sum a_i \beta_i$ » ; la fin de la ligne est biffée : « \struck{\ill{} pourvu que} »} \uncertain{lorsque} \uncertain{positifs} \struck{\ill{} \uncertain{soient} \ill{}} \note{en interligne : « \uncertain{on} \uncertain{peut} \uncertain{aussi} »} \ill{} \struck{\ill{} \uncertain{supposons} \uncertain{de} \uncertain{plus}, \uncertain{les} $a_i$} \note{ligne biffée} \uncertain{positifs} \struck{\ill{} \uncertain{dans} \uncertain{le} \uncertain{cas} b) \uncertain{on} \uncertain{a}} \note{« hypothèse » ou « suppose » en interligne} \uncertain{suppose} $\sum \alpha_i = \sum \beta_i$. \uncertain{On} \uncertain{a} \ill{} --- \uncertain{en} \uncertain{effet}
\begin{align}
  \text{(9)}\qquad \sum a_i \alpha_i &= (a_1 - a_2)\,\alpha_1 + (a_2 - a_3)(\alpha_1 + \alpha_2) + \cdots \notag \\
  &\qquad + (a_{n-1} - a_n)(\alpha_1 + \cdots + \alpha_{n-1}) + a_n(\alpha_1 + \cdots + \alpha_n) \notag
\end{align}
\note{les $\alpha$ de (9) sont surchargés, écrits sur des $\beta$ ou l'inverse ; le « (9) » est cerclé} \struck{\ill{}} \uncertain{et} \uncertain{une} \struck{$\sum a_i \ldots$} \uncertain{formule} \uncertain{analogue} \uncertain{pour} $\sum a_i \beta_i$. \uncertain{Comme} \uncertain{les} \uncertain{coefficients} $(a_1 - a_2), \ldots, (a_{n-1} - a_{n-1})$ \note{ainsi sur la page ; lire $(a_{n-1} - a_n)$} \uncertain{sont} \uncertain{positifs}, \struck{\ill{} \uncertain{ainsi} \uncertain{que} \ill{}} \uncertain{et} \uncertain{que} $\alpha_1 \leq \beta_1$, \ldots, $(\alpha_1 + \cdots + \alpha_{n-1}) \leq (\beta_1 + \cdots + \beta_{n-1})$, \uncertain{on} \uncertain{en} \uncertain{conclut} \uncertain{que} \uncertain{les} \uncertain{termes} \uncertain{de} \uncertain{la} \uncertain{somme} \ill{} (9) \uncertain{sont} \uncertain{majorés} \uncertain{par} \uncertain{les} \uncertain{termes} \uncertain{correspondants} \uncertain{de} \uncertain{la} \uncertain{somme} \uncertain{relative} \uncertain{aux} $\beta_i$, \uncertain{à} \uncertain{l'exception} \struck{\ill{}} \note{en interligne : « à priori »} \uncertain{du} \uncertain{dernier}. \uncertain{Mais} \uncertain{on} \uncertain{a} \struck{\ill{} \ill{}} \note{en interligne} \uncertain{aussi} $a_n(\alpha_1 + \cdots + \alpha_n) \leq a_n(\beta_1 + \cdots + \beta_n)$ \uncertain{dans} \uncertain{le} \uncertain{cas} \uncertain{a}) ($a_n \geq 0$, $\sum \alpha_i \leq \sum \beta_i$) \uncertain{et} \uncertain{dans} \uncertain{le} \uncertain{cas} \uncertain{b}) ($\sum \alpha_i = \sum \beta_i$), \uncertain{d'où} \struck{\uncertain{la} \uncertain{conclusion}} \note{en interligne : « l'inégalité (1) »}.

\emph{Corollaire \uncertain{1a}} \note{en interligne au-dessus, souligné : « Th. 1, 2°, b) » ; le numéro du corollaire est corrigé, « 1 » surchargé ; deux lignes sous le titre sont biffées, dont « \struck{(resp. \uncertain{du} \uncertain{corollaire} 1, a))} » et « \struck{les \uncertain{mêmes} \uncertain{hypothèses}} »} \uncertain{Sous} \uncertain{les} \uncertain{conditions} \uncertain{du} \struck{\ill{} \ill{}} \uncertain{Th}. 1, 2°, b), \struck{\ill{} \ill{}} \uncertain{on} \uncertain{a} \uncertain{aussi}
\[
  \text{(10)}\qquad f\Bigl(\frac{1}{\rho_1}, \ldots, \frac{1}{\rho_n}\Bigr) \leq f\Bigl(\frac{1}{\sigma_1}, \ldots, \frac{1}{\sigma_n}\Bigr)
\]
\note{le « (10) » est cerclé ; en marge à gauche, cerclé : « \uncertain{petits} \uncertain{caractères} », consigne de mise en page comme à la page 88}
\uncertain{En} \uncertain{effet}, \uncertain{on} \uncertain{a} $\rho_1 \cdots \rho_n = \sigma_1 \cdots \sigma_n$ \ill{} \uncertain{d'où} \uncertain{pour} \uncertain{tout} \ill{} \note{en marge : « $1 \leq M \leq n$ » ou « $\sigma \leq M \leq$ », incertain}
\[
  \rho_M \cdots \rho_n = \frac{\sigma_1 \cdots \sigma_{M-1}}{\rho_1 \cdots \rho_{M-1}}\; \sigma_M \cdots \sigma_n ,
  \qquad \text{d'où}\quad \rho_M \cdots \rho_n \;\uncertain{\geq}\; \sigma_M \cdots \sigma_n
\]
\note{« $\rho_M \cdots \rho_n$ » à gauche est surchargé, un « $|M|$ » ou « $\rho_M$ » repassé ; le facteur $\sigma_M \cdots \sigma_n$ est écrit après un mot biffé ; le sens des deux signes d'inégalité de cette ligne et de la suivante est peu lisible, le trait étant le même dans les deux cas} \ill{} $\sigma_1 \cdots \sigma_{M-1} \;\uncertain{\geq}\; \rho_1 \cdots \rho_{M-1}$. \uncertain{Il} \uncertain{s'ensuit} \uncertain{que} \uncertain{les} \uncertain{suites} $\bigl(\frac{1}{\rho_n}, \ldots, \frac{1}{\rho_1}\bigr)$ \uncertain{et} $\bigl(\frac{1}{\sigma_n}, \ldots, \frac{1}{\sigma_1}\bigr)$ \uncertain{satisfont} \uncertain{encore} \uncertain{aux} \uncertain{conditions} \uncertain{du} \uncertain{corollaire} (\uncertain{la} \uncertain{première} \uncertain{étant} \uncertain{décroissante}, \struck{\ill{} \ill{}} \uncertain{la} \uncertain{seconde} \ill{} \ill{} \uncertain{le} \uncertain{produit} \uncertain{des} $\frac{1}{\sigma_i}$, \uncertain{et} \uncertain{enfin} $\frac{1}{\rho_n} \cdots \frac{1}{\rho_i} \leq \frac{1}{\sigma_n} \cdots \frac{1}{\sigma_i}$ \uncertain{pour} $1 \leq M \leq n$), \uncertain{d'où} \uncertain{l'inégalité} (10).

\page{85}
\struck{\uncertain{Bien} \ill{} \uncertain{qu'on} \ill{} \uncertain{les} \uncertain{conditions} \ill{}} \struck{\uncertain{Bien} \ill{} \uncertain{les} \uncertain{conditions} \uncertain{du} \uncertain{Th}. 1, 1°, \uncertain{dans} \uncertain{les} \uncertain{hypothèses} \uncertain{b}), \ill{} \uncertain{que} \ill{} \ill{}} \struck{\uncertain{l'inégalité} \uncertain{analogue}} \struck{(\uncertain{11}) \qquad \ill{} $\varphi(-\alpha_1, \ldots, -\alpha_n) \leq \varphi(-\beta_1, \ldots, -\beta_n)$} \note{quatre lignes en tête de page, biffées d'un long trait ondulé et de traits verticaux ; la formule, numérotée « (11) » ou « (4) », est aussi surchargée}

\emph{Remarque} \note{souligné ; le mot recouvre un « \struck{Notons} » et un « \struck{toutefois} » biffés} \uncertain{que} \uncertain{si} \uncertain{deux} \uncertain{suites} $(\alpha_i)_{i \leq n}$, $(\beta_i)_{i \leq n}$ \uncertain{satisfont} \uncertain{aux} \note{en interligne : « ($\alpha_1 \geq \cdots \geq \alpha_n$, $\alpha_1 + \cdots + \alpha_m \leq \beta_1 + \cdots + \beta_m$ \ill{} $1 \leq m \leq n$) », cerclé} \uncertain{conditions} \uncertain{générales} \uncertain{du} \struck{\ill{} \uncertain{lemme}} \note{en interligne : « Th. 1, 1° »}, \struck{\ill{} \ill{} \uncertain{il} \ill{} \ill{} \uncertain{les} \ill{} \ill{}} \struck{\ill{} \uncertain{dans} \uncertain{la} \uncertain{suite} $\varphi(\alpha_i)_{i \leq p} \leq \varphi(\beta_i)_{i \leq p}$} \struck{(\uncertain{pour} \ill{} \ill{} $\sum \alpha_i \leq \sum \beta_i$ \ill{}} \struck{\uncertain{donc} \ill{} \ill{} \uncertain{pour} \ill{} \uncertain{alors} \uncertain{on} \uncertain{aura} \uncertain{aussi}} \struck{\ill{} \uncertain{que} \uncertain{les} \uncertain{fonctions} \uncertain{de} \ill{}} \struck{$\varphi(\alpha_1, \ldots, \alpha_p)$ \ill{} \uncertain{les} \ill{} \uncertain{pour}} \struck{\ill{} \uncertain{sous} \uncertain{les} \uncertain{conditions} \ill{} \ill{} \ill{} $\alpha_{p+1}, \ldots, \alpha_n$ \uncertain{et} $\beta_{p+1}, \ldots, \beta_n$} \note{bloc de six lignes biffé de grands traits ondulés, dont on ne retient que les formules} \uncertain{alors} \uncertain{il} \uncertain{en} \uncertain{est} \uncertain{de} \uncertain{même} \note{au bas du bloc, souligné : « à la définition »} \uncertain{des} \uncertain{suites} $(\alpha_i)_{i \leq k}$ \uncertain{et} $(\beta_i)_{i \leq k}$ \uncertain{pour} \uncertain{tout} \uncertain{entier} $k$ \uncertain{tel} \uncertain{que} \note{en interligne : « qui suit »} $1 \leq k \leq n$. \uncertain{Comme} \ill{} \ill{} \uncertain{à} \uncertain{une} \uncertain{suite} \uncertain{de} \ill{} \uncertain{le} \uncertain{lemme}, \uncertain{et} \uncertain{que} \uncertain{si} $\varphi$ \uncertain{est} \uncertain{une} \uncertain{fonction} \uncertain{convexe} \struck{\uncertain{symétrique} \uncertain{sur} \ill{} \uncertain{les} \ill{}} \note{en interligne : « croissante » ou « \uncertain{de} $k$ »} \uncertain{de} $k$ \uncertain{arguments}, \uncertain{on} \uncertain{a} \uncertain{aussi}
\[
  \varphi(\alpha_1, \ldots, \alpha_k) \leq \varphi(\beta_1, \ldots, \beta_k) .
\]
\uncertain{La} \uncertain{même} \uncertain{remarque} \ill{} \uncertain{pour} \ill{} 2 \ill{} \uncertain{conditions} \uncertain{des} \struck{\uncertain{corollaire} 1 \uncertain{du} \uncertain{lemme}} \note{en interligne, remplaçant le texte biffé : « Th. 1, 2° »}.

\struck{\uncertain{Applications} \uncertain{des} \uncertain{corollaires} 1} \note{ligne biffée, soulignée d'un trait double} \uncertain{Donnons} \uncertain{des} \uncertain{cas} \uncertain{intéressants} \uncertain{de} \uncertain{fonctions} $f$ \uncertain{satisfaisant} \uncertain{aux} \uncertain{conditions} \uncertain{du} \struck{\uncertain{corollaire} 1} \note{en interligne : « Th. 1, 2° »}. \struck{\ill{}} \uncertain{Si} \struck{\ill{}} $f(\rho_1, \ldots, \rho_n)$ \uncertain{est} \uncertain{une} \uncertain{fonction} \uncertain{convexe} \uncertain{croissante} \note{souligné} \uncertain{des} \uncertain{arguments} $\rho_i > 0$, \uncertain{alors} $f(e^{t_1}, \ldots, e^{t_n})$ \uncertain{est} \struck{\ill{}} \note{en interligne : « une »} \struck{\ill{} \ill{} \ill{} \uncertain{les} \uncertain{arguments}} $(t_i) \in \mathbf{R}^n$ \struck{\ill{} \uncertain{fonction} \uncertain{convexe} \ill{} \uncertain{qui} \ill{}} \note{en interligne, souligné : « $f$ sera » ; sept lignes biffées, avec en interligne « \uncertain{limite} \uncertain{de} \uncertain{fonctions} \uncertain{linéaires} »} $f(e^{t_1}, \ldots, e^{t_n})$ \uncertain{est} \uncertain{enveloppe} \uncertain{supérieure} \struck{\ill{} \uncertain{fonction} \uncertain{convexe}} \uncertain{des} \uncertain{fonctions} $\sum a_i e^{t_i} + b$, \uncertain{qui} \uncertain{sont} \uncertain{convexes} \uncertain{puisque} \uncertain{les} \uncertain{fonctions} $e^{t_i}$ \uncertain{sont} \uncertain{convexes}, \uncertain{et} \uncertain{les} $a_i$ \ill{}, \uncertain{elle} \uncertain{est} \uncertain{elle-même} \uncertain{convexe}. \uncertain{Enfin} \struck{\ill{} \uncertain{pour} \ill{} \uncertain{une} \uncertain{fonction} \uncertain{de} \ill{}} \struck{\ill{}, \uncertain{comme} \uncertain{une} \uncertain{de} \ill{} \uncertain{sur} $\mathbf{R}^n$ \uncertain{toute}} \note{deux lignes biffées, soulignées} \emph{\uncertain{norme}} \note{souligné d'un trait double, avec un trait rouge dessous} $N(x) = N(x_1, \ldots, x_n)$ \uncertain{qui} \uncertain{est} \uncertain{symétrique}, \uncertain{telle} \uncertain{que} \note{en interligne : « pour $x_1 \geq 0, \ldots, x_n \geq 0$ »}
\[
  \text{(\uncertain{12})}\qquad N(x_1, \ldots, x_n) = N(|x_1|, \ldots, |x_n|) \qquad \struck{\ill{} \uncertain{fonction} \uncertain{croissante} \uncertain{de} \uncertain{chaque} \ill{}}
\]
\note{le numéro est cerclé ; la fin de la ligne est biffée, ainsi que les deux lignes suivantes : « \struck{\ill{} \uncertain{normalisée} \uncertain{par} $N(e_i) = 1$} », « \struck{\ill{} \uncertain{les} \ill{} \uncertain{sont} \ill{}} »}
\marginal{(\uncertain{Nous} \ill{} \uncertain{supposons} $N$ \ill{} \ill{} \uncertain{vérifiant} \uncertain{les} \uncertain{conditions} $(\lambda_i)$ \ill{} \ill{} \uncertain{des} \ill{} (\uncertain{i.e.} 2 \uncertain{les} \uncertain{conditions}, \uncertain{et} \ill{} \uncertain{stable} \ill{} \uncertain{multiplication} \uncertain{par} \ill{} $(\lambda_i)$ \uncertain{avec} $|\lambda_i| = 1$ \ill{} \ill{}) \ill{} \ill{} \uncertain{majorant} \ill{} \ill{} \ill{} $\sum \lambda_i$ \ill{} \ill{} $f$ \uncertain{tel} \ill{} $\sum d_i \leq \sum$ \ill{} \ill{} \uncertain{seulement} \ill{} \ill{} $n$ \uncertain{éléments} \ill{} \ill{})} \note{longue note de côté dans la marge gauche, en deux blocs, l'un descendant le long de la page et l'autre en escalier au bas ; un trait rouge la longe ; elle porte sur la normalisation de $N$ et reste pour l'essentiel illisible même redressée}

\page{86}
\uncertain{norme} \uncertain{de} \uncertain{Schatten} \note{ainsi lu ; le mot est abrégé « Sch. » et suivi d'un trait} $N(x_1, \ldots, x_n)$ \uncertain{satisfait} \uncertain{aux} \uncertain{conditions} \uncertain{exigées} \note{en interligne : « id \ill{} »} dans l'énoncé du \struck{\uncertain{corollaire}} \note{en interligne : « Th. 1, 2° »} \note{en interligne : « pour la fonction $f$. On en \ill{} \uncertain{déduit} »} \struck{\ill{} \uncertain{application} \ill{} \uncertain{suivante}} : \struck{si $(\rho_i)$ \struck{$\rho(x)$} et $(\sigma_i)$ \uncertain{sont} \uncertain{comme} \uncertain{dans} \uncertain{le} \uncertain{Th}.} \struck{\ill{} $=$ \ill{} \uncertain{soit} \ill{} \uncertain{la} \uncertain{fonction}} \struck{\ill{} \ill{}} \struck{\ill{} \uncertain{pour} \uncertain{tout} \ill{}} \struck{\ill{} \uncertain{conditions}, \ill{} \ill{}} \struck{\uncertain{dans} \uncertain{l'hypothèse} \uncertain{b}), \ill{}} \struck{\uncertain{dans} \uncertain{l'hypothèse} a) \ill{}} \note{bloc de sept lignes entouré d'un cadre et biffé ; sa fin est reprise en clair :} \uncertain{dans} \uncertain{l'hypothèse} \uncertain{b}), \ill{} $\pm 0$) \ill{} \uncertain{aussi}. \uncertain{Que} \uncertain{la} \uncertain{fonction} $(e^{t_i})^p$ \uncertain{est} \ill{} \uncertain{aussi} $f(r_1, \ldots, r_n) = r_1^p \cdots r_n^p$ \note{formule en partie encadrée d'un trait ondulé ; l'exposant $p$ du dernier facteur est incertain} \ill{} \uncertain{qu'un} $p$ \uncertain{réel} \uncertain{quelconque} \uncertain{est} \uncertain{satisfait} \uncertain{aux} \ill{} \ill{} \ill{} \uncertain{conditions} \struck{\ill{}} \uncertain{dans} \uncertain{le} \uncertain{corollaire} 1, \struck{\ill{}} \uncertain{et} $=$ \struck{\ill{} \ill{}} \uncertain{conditions} \uncertain{ainsi} \uncertain{que} \ill{} \uncertain{Pour} \uncertain{tout}, \ill{} \uncertain{les} \uncertain{conditions} \ill{} \note{en interligne : « \ill{} \uncertain{réel} \ill{} »} \uncertain{b}), \ill{} ($\prod_1^m \rho_i \leq \prod_1^m \sigma_i$, $\prod_1^n \rho_i = \prod_1^n \sigma_i$) \note{« $m$ » et « $n$ » au-dessus des $\prod$ ; le second $\prod$ est corrigé}
\[
  \text{(15)}\qquad \sum_1^n \rho_i^p \leq \sum_1^n \sigma_i^p \qquad (p \text{ réel})
\]
\note{le « (15) » est surchargé, écrit sur un « (9) » ou un « (1) »} \struck{\ill{} \ill{} \ill{}, \ill{}} \uncertain{et} \uncertain{si} $p \geq 0$, \struck{\ill{}} \note{en interligne : « $\prod_1^m \rho_i \leq \prod_1^m \sigma_i$ \ill{} »}
\[
  \text{(16)}\qquad \sum_1^k \rho_i^p \leq \sum_1^k \sigma_i^p \qquad (p \geq 0,\; \struck{\ill{}} k = 1, \ldots, n) .
\]
\note{le « (16) » est écrit au-dessus d'un « (12) » biffé ; un mot est biffé après « $p \geq 0$, »}
\struck{\uncertain{Remarquons} \uncertain{d'ailleurs} \uncertain{que} \uncertain{deux} \uncertain{systèmes} $(\rho_i)$, $(\sigma_i)$ \ill{} \uncertain{satisfont} \uncertain{aux} \ill{} \uncertain{inégalités} (16), \ill{} \uncertain{pour} \ill{} $p$ \uncertain{entier} $\geq 0$, \ill{} \uncertain{pour} \ill{} \uncertain{vérifient} \ill{}} \note{quatre lignes entourées d'un cadre et barrées de grands zigzags}

\struck{\ill{} \ill{}}

\emph{Th 2} (\uncertain{Premier} \uncertain{théorème} \uncertain{de} \uncertain{convexité} \uncertain{dans} \uncertain{les} \uncertain{espaces} \uncertain{de} \uncertain{Hilbert}) \note{« Th 2 » dans la marge gauche, souligné ; le titre entre parenthèses est souligné}

\uncertain{Soit} $u$ \uncertain{un} \uncertain{opérateur} \struck{\ill{}} \uncertain{dans} \uncertain{un} \uncertain{espace} \uncertain{de} \uncertain{Hilbert}, \uncertain{soit} \struck{$\varphi$ \uncertain{la} \uncertain{suite}} $f(\rho_1, \ldots, \rho_n)$ \note{en interligne} \uncertain{une} \uncertain{fonction} (\uncertain{de} $n$ \uncertain{arguments} $r_i \geq 0$, \uncertain{croissante} \uncertain{pour} \uncertain{chaque} \uncertain{argument}, \uncertain{et} \uncertain{telle} \uncertain{que} $f(e^{t_1}, \ldots, e^{t_n})$ \uncertain{soit} \uncertain{fonction} \uncertain{convexe} \uncertain{de} $(t_i) \in \mathbf{R}^n$. \uncertain{Alors} \uncertain{on} \uncertain{a} \struck{\ill{}} \note{« $N_2$ » ou « $N_0$ » biffé, puis en interligne : « (\uncertain{avec} \uncertain{les} \uncertain{notations} \uncertain{ordinaires} \uncertain{du} \uncertain{n°} 3) » ; l'énoncé est marqué d'un trait vertical dans la marge}
\[
  \text{(17)}\qquad f(|\lambda_1|, \ldots, |\lambda_n|) \leq f(\rho_1, \ldots, \rho_n)
\]
\uncertain{Il} \uncertain{suffit} \uncertain{en} \uncertain{effet} \uncertain{d'appliquer} \uncertain{le} \uncertain{corollaire} 1 \uncertain{du} \struck{\uncertain{lemme}} \note{en interligne : « Th. 1, 2° »} (\uncertain{hypothèse} a)) \uncertain{avec} \ill{} $N_2$ \note{ainsi ; peut-être « $\rho_i$ » ou le n° 2}, \uncertain{pour} 3 ; \uncertain{corollaire} 2. --- \uncertain{En} \uncertain{particulier} \ill{} \note{ligne obscure : elle renvoie aux corollaires du n° 3 pour les inégalités $|\lambda_1 \cdots \lambda_k| \leq \rho_1 \cdots \rho_k$}

\emph{Corollaire 1.} Si $f(\rho)$ \uncertain{est} \uncertain{une} \uncertain{fonction} \struck{\uncertain{convexe}} \note{en interligne : « croissante »} \uncertain{croissante} \uncertain{de} $\rho \geq 0$ \uncertain{telle} \uncertain{que} $f(e^t)$ \uncertain{soit} \uncertain{convexe}, \uncertain{et} $u$ \uncertain{un} \uncertain{opérateur} \uncertain{dans} \uncertain{un} \uncertain{espace} \uncertain{de} \uncertain{Hilbert}, \uncertain{on} \uncertain{a}
\marginal{(13) $N(\rho_1, \ldots, \rho_n) \leq N(\sigma_1, \ldots, \sigma_n)$ \ill{} \ill{} \ill{} \ill{} \uncertain{lorsque} $N$ \ill{} \ill{} \uncertain{pour} \uncertain{tout} \ill{} $1 \leq k \leq n$ (14) $N(\rho_1, \ldots, \rho_k) \leq N(\sigma_1, \ldots, \sigma_k)$, $1 \leq k \leq n$ \ill{} \ill{} \uncertain{pour} \ill{}} \note{note de côté dans la marge gauche, en regard du haut de la page : deux formules cerclées (13) et (14), la seconde écrite deux fois, l'une biffée, et quelques mots de commentaire illisibles ; les numéros sont lus (13) et (14) par leur place entre (12) et (15), mais le second chiffre est peu net}

\page{87}
\note{Petit feuillet brunâtre, comme les pages 74 et 75 du lot 4, écrit sur sa moitié supérieure seulement : une reprise au propre de la fin de la page 85 et du haut de la page 86, avec ses propres numéros (13) et (14).}
\uncertain{En} \uncertain{anticipant} \uncertain{l'un}, \uncertain{soit} $f(\rho)$ \uncertain{une} \uncertain{fonction} \uncertain{définie} \uncertain{pour} $\rho > 0$, \uncertain{telle} \uncertain{que} $f(e^t)$ \uncertain{soit} \uncertain{convexe}, \uncertain{posons} $f(\rho_1, \ldots, \rho_n) = \sum f(\rho_i)$, \uncertain{c'est} \uncertain{une} \uncertain{fonction} \uncertain{symétrique} \uncertain{des} \note{en interligne : « $\rho_i$ »} \struck{\ill{}} \uncertain{des} $\rho_i > 0$, \uncertain{et} \uncertain{telle} \uncertain{que} $f(e^{t_1}, \ldots, e^{t_n})$ \uncertain{soit} \struck{\ill{}} \uncertain{convexe}. \uncertain{Par} \uncertain{suite}, \struck{\ill{}} $f(\rho_1, \ldots, \rho_n)$ \uncertain{satisfait} \uncertain{aux} \uncertain{conditions} \uncertain{du} \uncertain{corollaire} 1, \uncertain{sous} \uncertain{l'hypothèse} \uncertain{b}. \uncertain{Si} \uncertain{de} \uncertain{plus} $f(\rho)$ \uncertain{est} \uncertain{croissante}, \uncertain{et} \uncertain{définie} \uncertain{et} \uncertain{continue} \uncertain{aussi} \uncertain{pour} $\rho \geq 0$, \uncertain{alors} $f(\rho_1, \ldots, \rho_n)$ \uncertain{satisfait} \uncertain{aux} \uncertain{conditions} \uncertain{du} \struck{\uncertain{corollaire} 1} \note{en interligne : « Th. 1, 2° »} \uncertain{sous} \uncertain{l'hypothèse} a. \uncertain{D'où}, \struck{\ill{} \ill{} \ill{}} \note{en interligne : « \uncertain{sous} \uncertain{les} \uncertain{conditions} a), b), \uncertain{respectivement} »} \struck{$(\alpha_i)$, $(\beta_i)$, \ill{} \uncertain{inégalités}} \uncertain{l'inégalité} :
\[
  \text{(13)}\qquad \sum_1^n f(\rho_i) \leq \sum_1^n f(\sigma_i)
\]
\uncertain{et} \uncertain{sous} \uncertain{les} \uncertain{conditions} a), \uncertain{compte} \uncertain{tenu} \uncertain{de} \uncertain{la} \uncertain{remarque} \uncertain{précédente}, \uncertain{la} \uncertain{suite} \uncertain{d'inégalités}
\[
  \text{(14)}\qquad \sum_1^k f(\rho_i) \leq \sum_1^k f(\sigma_i) \qquad (k = 1, \ldots, n)
\]
\uncertain{Prenons} \uncertain{p.ex}, \struck{\uncertain{pour} \ill{}} \note{en interligne : « \uncertain{pour} $p$ \uncertain{réel} \uncertain{positif} »} $f(\rho) = \rho^p$, \note{en interligne : « ($\rho > 0$) »} \uncertain{alors} $f(e^t) = e^{pt}$ \note{« $f(\rho)$ » écrit sous $f(e^t)$}, \uncertain{qui} \uncertain{est} \uncertain{bien} \uncertain{une} \uncertain{fonction} \uncertain{convexe}, \struck{\ill{}} \uncertain{et} \uncertain{définie} \uncertain{et} \struck{\ill{}} \uncertain{continue} \struck{\ill{}} \uncertain{fonction} \uncertain{croissante} \uncertain{de} $\rho$, \uncertain{pour} $\rho \geq 0$, \uncertain{et} \struck{\ill{}} \uncertain{d'où} \uncertain{pour} $p \geq 0$ \uncertain{les} \struck{\ill{} \ill{}} \note{la dernière ligne est biffée ; le reste de la feuille est blanc}

\page{88}
\[
  \text{(18)}\qquad \struck{f(\rho)}\; \sum_1^k f(|\lambda_i|) \leq \sum_1^k f(\rho_i) \qquad (k = 1, \ldots) \; \struck{(k \leq \dim E)}
\]
\note{les formules (18) et (19) sont marquées d'un trait vertical dans la marge}
\uncertain{En} \uncertain{particulier}, \struck{\uncertain{pour} \ill{}} \uncertain{si} \ill{} \uncertain{est} \uncertain{de} \uncertain{dimension} \ill{}, \uncertain{si} $\sum_{i=0}^\infty f(\rho_i)$ \struck{\uncertain{converge}, \ill{}} \note{en interligne : « \uncertain{est} \uncertain{convergente} »}, \uncertain{il} \uncertain{en} \uncertain{est} \uncertain{de} \uncertain{même} \uncertain{de} $\sum_{i=0}^\infty f(|\lambda_i|)$, \uncertain{et} \uncertain{on} \uncertain{a}
\[
  \text{(19)}\qquad \sum_{i=1}^\infty f(|\lambda_i|) \leq \sum_{i=0}^\infty f(\rho_i)
\]
\note{ainsi sur la page : la somme de gauche part de $i = 1$, celle de droite de $i = 0$}

\uncertain{En} \uncertain{particulier}, \uncertain{prenons} $f(r) = r^p$ ($p > 0$), \ill{} \uncertain{telle} \ill{} \uncertain{que} \ill{} \uncertain{compris}, \uncertain{d'où} \ill{} \ill{} \uncertain{le} \ill{} $p > 0$. \uncertain{Alors}, \uncertain{pour} \uncertain{tout} \uncertain{entier} $k > 0$

\emph{Corollaire 2} \uncertain{Soit} $u$ \uncertain{un} \uncertain{opérateur} \uncertain{dans} \uncertain{un} \uncertain{espace} \uncertain{de} \uncertain{Hilbert} $E$, \uncertain{et} \uncertain{soit} $p > 0$ \uncertain{un} \uncertain{nombre} \uncertain{réel}. \uncertain{Alors} \uncertain{on} \uncertain{a} \note{l'énoncé est marqué d'un trait vertical dans la marge}
\[
  \text{(20)}\qquad \sum_{i=1}^k |\lambda_i|^p \leq \sum_{i=1}^k \rho_i^p = \struck{\|u\|_p} \qquad \struck{(k \leq \dim E)}
\]
\note{formule encadrée au crayon rouge ; le « $= \|u\|_p$ » final est cerclé et biffé, la parenthèse de droite est barrée d'un zigzag ; sous la formule, une ligne biffée : « \struck{\ill{} $E$ \uncertain{est} \uncertain{de} \uncertain{dimension} \uncertain{finie} \ill{}} »}
\[
  \text{(21)}\qquad \sum_i |\lambda_i|^p \leq \sum_i \rho_i^p = \|u\|_p
\]
\note{formule encadrée au crayon rouge} (\uncertain{où} \uncertain{le} \uncertain{second} \uncertain{membre} \uncertain{est} \uncertain{fini} \uncertain{ou} \uncertain{non}, \uncertain{la} \uncertain{somme} \uncertain{est} \ill{})

\uncertain{En} \uncertain{particulier}, \uncertain{prenons} $p = 1$ \uncertain{et} $p = 2$, \uncertain{et} \uncertain{utilisant} \uncertain{le} \uncertain{N°} 1, \uncertain{Th}. 2 \uncertain{et} \uncertain{Th}. 3, \uncertain{formules} (15) \uncertain{et} (20), \uncertain{on} \uncertain{trouve}

\emph{Corollaire 3} \uncertain{Soit} $u$ \uncertain{un} \uncertain{opérateur} \uncertain{de} \uncertain{Fredholm} (\uncertain{resp}. \uncertain{un} \uncertain{op}. \uncertain{de} \struck{\uncertain{Hilb}} \uncertain{H}.-\uncertain{S}.) \uncertain{dans} \struck{\ill{}} \uncertain{un} \uncertain{espace} \uncertain{de} \uncertain{Hilbert}. \uncertain{Alors} \uncertain{la} \uncertain{suite} \uncertain{des} \uncertain{valeurs} \uncertain{propres} \uncertain{est} \uncertain{sommable} (\uncertain{resp}. \uncertain{de} \uncertain{carré} \uncertain{sommable}) \uncertain{et} \uncertain{on} \uncertain{a}
\[
  \text{(21)}\qquad \sum |\lambda_i| \leq \|u\|_1
\]
\[
  \text{(22)}\qquad \Bigl(\sum |\lambda_i|^2\Bigr)^{1/2} \leq \|u\|_2
\]
\note{les deux formules sont encadrées ensemble au crayon rouge ; la première porte le numéro (21) comme la formule précédente, ainsi sur la page}

\uncertain{Signalons} \uncertain{aussi} \uncertain{le}

\emph{Corollaire 4} \uncertain{Soit} $u$ \uncertain{un} \uncertain{opérateur} \struck{\ill{}} \uncertain{dans} \uncertain{un} \uncertain{espace} \uncertain{de} \uncertain{Hilbert} \uncertain{de} \uncertain{dimension} \uncertain{finie} $n$, \struck{\uncertain{Alors}} \uncertain{soit} $f(\rho_1, \ldots, \rho_n)$ \uncertain{une} \uncertain{fonction} \uncertain{de} $n$ \uncertain{arguments} $\rho_i > 0$ \uncertain{telle} \uncertain{que} $f(e^{t_1}, \ldots, e^{t_n})$ \uncertain{soit} \uncertain{convexe} (\uncertain{pas} \uncertain{forcément} \uncertain{croissante}). \uncertain{Si} $u$ \uncertain{est} \uncertain{défini} \uncertain{pour} \struck{\ill{}} $\rho_i \geq 0$, \uncertain{ou} \uncertain{si} $u$ \uncertain{est} \uncertain{inversible}, \uncertain{on} \uncertain{a} \note{l'énoncé est marqué d'un trait vertical dans la marge, et d'un trait rouge}
\[
  \text{(23)}\qquad f(|\lambda_1|, \ldots, |\lambda_n|) \leq f(\rho_1, \ldots, \rho_n) .
\]
\note{le « (23) » est cerclé ; la formule est encadrée}
\marginal{\uncertain{petits} \uncertain{caractères}} \note{écrit de côté dans la marge gauche, en regard du corollaire 4 : une consigne de mise en page, non une note mathématique}
\struck{\uncertain{Alors}} \uncertain{En} \uncertain{effet}, \ill{} \uncertain{résulte} \uncertain{des} \uncertain{inégalités}
\[
  |\lambda_1 \cdots \lambda_k| \leq \rho_1 \cdots \rho_k \qquad (k \leq n),
\]
\uncertain{l'égalité}

\page{89}
$|\lambda_1| \cdots |\lambda_n| = \rho_1 \cdots \rho_n$. \uncertain{En} \uncertain{effet}, \struck{\ill{}} \note{ce qui suit est écrit dans deux cadres emboîtés}
\[
  |\lambda_1 \cdots \lambda_n|^2 = |\det u|^2 = \ill{} = \det u^{*} u = \rho_1^2 \cdots \rho_n^2 \quad (\text{\uncertain{de} } u^{*} u)
\]
\ill{} \uncertain{les} $\rho_i^2$ \uncertain{étant} \uncertain{les} \uncertain{valeurs} \uncertain{propres} \uncertain{de} $u^{*} u$ \ill{}, \uncertain{d'où} \ill{} \uncertain{i.e.} \ill{} \uncertain{la} \uncertain{condition} \ill{} \uncertain{du} \uncertain{Th}. 1, 2° \ill{} \uncertain{b}) \ill{} \uncertain{vérifiée}. \ill{} \uncertain{les} \uncertain{conditions} \ill{} \note{trois lignes sous les cadres, très rapides}
\marginal{\struck{(25) $\lambda_1(u) \cdots \lambda_n(u) = \det u$} $\bigl(\rho_1(u) \cdots \rho_n(u)\bigr)^2 = \rho_1(u^{*}u) \cdots \rho_n(u^{*}u) = \det(u^{*}u) = \det u^{*} \det u = \overline{\det u}\, \det u = |\det u|^2$, d'où (26) $\rho_1(u) \cdots \rho_n(u) = |\det u|$} \note{de côté dans la marge gauche, en haut de la page ; la formule (26), cerclée, est encadrée ; le numéro (25) est biffé et remplacé ; un « (24) » figure plus bas dans la même marge, près d'un bloc au crayon gribouillé jusqu'à l'illisibilité}
\uncertain{Si} $u$ \uncertain{est} \uncertain{inversible}, \ill{} \uncertain{dans} \ill{} \uncertain{les} \ill{} \uncertain{les} \uncertain{conditions} \uncertain{du} \uncertain{Th}. 1, 2°, \uncertain{hyp}. \uncertain{b}) \uncertain{sont} \uncertain{vérifiées}, \ill{} \ill{}. \uncertain{Si} $f$ \uncertain{est} \ill{} \uncertain{définie} \uncertain{et} \uncertain{continue} \uncertain{pour} $\rho_i \geq 0$, \ill{} \uncertain{aussi} (23), \uncertain{pour} $u$ \uncertain{non} \uncertain{inversible}, \ill{} \uncertain{résulte} \uncertain{de} \ill{} \uncertain{continuité} \ill{} \uncertain{vérifiée} \uncertain{aussitôt}, \ill{} \uncertain{vérifiée} \uncertain{pour} \struck{\uncertain{tout}} $u$ \ill{}. \struck{\ill{} \ill{}} \struck{\ill{} \uncertain{corollaire} 4 \ill{} \uncertain{par} \ill{} \uncertain{inversible} \ill{} \uncertain{la} \uncertain{formule} (21) \ill{} \uncertain{vérifiée} \ill{}} \note{trois lignes biffées d'un trait ondulé} \struck{\uncertain{Nota} \ill{}} \struck{(26) \ill{}} \note{deux lignes biffées}

\emph{Corollaire 5} \uncertain{Soit} $u$ \uncertain{un} \uncertain{opérateur} \uncertain{dans} \uncertain{un} \uncertain{espace} \uncertain{de} \uncertain{Hilbert} $E$, \uncertain{soit} $N$ \uncertain{une} \uncertain{norme} \uncertain{de} \uncertain{Schatten} \uncertain{sur} $\mathbf{R}^k$ \note{en interligne : « $k \leq \dim E$ » ; un « $I = (1, \ldots, n)$ » biffé}, \ill{} \uncertain{alors} \note{l'énoncé est marqué d'un trait vertical dans la marge}
\[
  \text{(27)}\qquad N(|\lambda_1|, \ldots, |\lambda_k|) \leq N(\rho_1, \ldots, \rho_k)
\]
\struck{\uncertain{D'où}, \ill{}} \uncertain{Par} \uncertain{suite}, \uncertain{si} \uncertain{la} \uncertain{dimension} \uncertain{de} $E$ \uncertain{est} \uncertain{infinie}, \uncertain{et} \uncertain{si} $N$ \uncertain{est} \uncertain{une} \uncertain{norme} \uncertain{de} \uncertain{Schatten} \uncertain{sur} \uncertain{l'espace} $\mathbf{R}^{(\mathbf{N})}$, \struck{\ill{} $N$ \ill{}} \note{en interligne : « \ill{} »}
\[
  \text{(28)}\qquad N\bigl((|\lambda_i|)\bigr) \leq N\bigl((\rho_i)\bigr)
\]
\uncertain{Dans} \uncertain{cet} \uncertain{énoncé}, $\mathbf{N}$ \uncertain{désigne} \uncertain{l'ensemble} \uncertain{des} \uncertain{entiers}, \note{le $\mathbf{N}$ est souligné, et de même dans $\mathbf{R}^{(\mathbf{N})}$} $\mathbf{R}^{(\mathbf{N})}$ \uncertain{l'espace} \uncertain{des} \uncertain{suites} \uncertain{de} \uncertain{nombres} \uncertain{réels} \uncertain{qui} \uncertain{sont} \ill{} \uncertain{nulles} \ill{} (\uncertain{i.e.} \uncertain{n'ayant} \uncertain{qu'un} \uncertain{nombre} \uncertain{fini} \uncertain{de} \uncertain{termes} \uncertain{non} \uncertain{nuls}). \uncertain{Une} \uncertain{norme} \uncertain{sur} $\mathbf{R}^{(\mathbf{N})}$ \uncertain{est} \uncertain{dite} \uncertain{une} \uncertain{norme} \uncertain{de} \uncertain{Schatten}, \uncertain{si} \uncertain{elle} \uncertain{induit} \uncertain{une} \uncertain{norme} \uncertain{de} \uncertain{Schatten} \uncertain{sur} \uncertain{tout} $\mathbf{R}^n$, \ill{} $\mathbf{R}^n$ (\uncertain{identifié} \uncertain{à} \uncertain{un} \uncertain{sous-ensemble} \uncertain{de} $\mathbf{R}^{(\mathbf{N})}$).

\page{90}
\note{Un feuillet est collé sur le bas de la page et en couvre les dernières lignes ; il est transcrit à sa place, après le corollaire 1.}
\struck{\ill{} \ill{}} \note{deux mots biffés en tête de page} \uncertain{En} \uncertain{effet}, \uncertain{si} $\xi = (\xi_i)$ \uncertain{est} \uncertain{un} \uncertain{élément} \uncertain{de} $\mathbf{R}^{(\mathbf{N})}$ \ill{} \uncertain{positif}, \uncertain{posons} $x_n = (\xi_1, \ldots, \xi_n, 0, \ldots, 0)$, \uncertain{alors} $N(x_n)$ \struck{\uncertain{est} \uncertain{une} \uncertain{suite} \uncertain{croissante} \uncertain{et} \ill{}} \note{« \uncertain{tend} \uncertain{avec} \ill{} » en interligne} \uncertain{on} \uncertain{pose}
\[
  N(\xi) = N\bigl((\xi_i)\bigr) = \lim N(x_n), \qquad \ill{} \text{ limite croissante } \geq 0 .
\]
\note{« limite croissante » incertain}
\struck{\uncertain{Alors}} \uncertain{La} \uncertain{formule} (28, \ill{}) \ill{} \uncertain{particulier} \uncertain{de} \ill{}, \uncertain{si} $N$ \struck{\ill{} \ill{}} \note{en interligne : « \ill{} »} \ill{}, \ill{} $f$ \ill{}, \ill{}. \uncertain{La} \uncertain{formule} (28, \ill{}) \uncertain{en} \uncertain{résulte} \uncertain{par} \uncertain{un} \uncertain{passage} \uncertain{à} \uncertain{la} \uncertain{limite}.

\emph{Th 3} \emph{Deuxième théorème de convexité dans les espaces de Hilbert} \note{« Th 3 » dans la marge, souligné ; le titre est souligné}

\uncertain{Soient} $E, F, G$ \uncertain{trois} \uncertain{espaces} \uncertain{de} \uncertain{Hilbert}, \uncertain{soit} $u$ \uncertain{une} \uncertain{application} \uncertain{linéaire} \struck{\uncertain{continue}} \note{en interligne : « cpte »} \uncertain{de} $E$ \uncertain{dans} $F$, \struck{\uncertain{et}} $v$ \uncertain{une} \uncertain{application} \uncertain{linéaire} \struck{\uncertain{continue}} \note{en interligne : « cpte »} \uncertain{de} $F$ \uncertain{dans} $G$, \struck{\ill{} \ill{} \ill{}} \note{en interligne : « \uncertain{symétrique} \uncertain{et} \uncertain{croissante} » puis « soit »} \uncertain{soit} $f(\rho_1, \ldots, \rho_n)$ \uncertain{une} \uncertain{fonction} \struck{\ill{}} (\uncertain{de} $n$ \uncertain{arguments} $\rho_i \geq 0$, \uncertain{telle} \uncertain{que} \struck{\ill{}} $f(e^{t_1}, \ldots, e^{t_n})$ \uncertain{soit} \uncertain{fonction} \uncertain{convexe} \uncertain{de} $(t_i) \in \mathbf{R}^n$. \uncertain{Alors} \uncertain{on} \uncertain{a} \note{l'énoncé est marqué d'un trait vertical dans la marge}
\[
  \text{(29)}\qquad f\bigl(\rho_1(vu), \ldots, \rho_n(vu)\bigr) \leq f\bigl(\rho_1(v)\rho_1(u), \ldots, \rho_n(v)\rho_n(u)\bigr)
\]
\note{le « (29) » est cerclé et surchargé}
\uncertain{En} \uncertain{effet}, \uncertain{il} \uncertain{suffit} \uncertain{d'appliquer} \uncertain{le} \struck{\uncertain{lemme}} \note{en interligne : « th. 1 »} \uncertain{fondamental}, \struck{\ill{}} \uncertain{conditions} \uncertain{aux} \uncertain{deux} \uncertain{suites} $\bigl(\rho_i(vu)\bigr)_{i \leq n}$ \uncertain{et} $\bigl(\rho_i(v)\rho_i(u)\bigr)_{i \leq n}$. \uncertain{Les} \uncertain{conditions} \struck{\uncertain{du} \uncertain{corollaire} 1 \uncertain{de} \uncertain{ce} \uncertain{lemme} \uncertain{sont}} \note{en interligne : « à l'hyp. b », « n° 2°) du th. 1 »} \uncertain{satisfaites}, \struck{\ill{}} \uncertain{en} \uncertain{vertu} \uncertain{du} \uncertain{N°} 3, \uncertain{Prop}. 3, \uncertain{corollaire} 3 \struck{\ill{}}. \note{la fin de la ligne est biffée}

\emph{Corollaire 1} \uncertain{Soient} $u$ \uncertain{et} $v$ \uncertain{comme} \uncertain{dans} \uncertain{le} \uncertain{Th}. 3, \uncertain{soit} $f(\rho)$ \uncertain{une} \uncertain{fonction} \uncertain{continue} \uncertain{croissante} \uncertain{de} $\rho \geq 0$ \uncertain{telle} \uncertain{que} $f(e^t)$ \uncertain{soit} \struck{\ill{}} \uncertain{fonction} \uncertain{convexe} \uncertain{de} $t \in \mathbf{R}$. \uncertain{Alors}, \uncertain{pour} \struck{\uncertain{tout}} \struck{\ill{}} \uncertain{entier} $n > 0$, \uncertain{on} \uncertain{a} \note{l'énoncé est marqué d'un trait vertical dans la marge}
\[
  \text{(30)}\qquad \sum_{i=1}^n f\bigl(\rho_i(vu)\bigr) \leq \sum_{i=1}^n f\bigl(\rho_i(v)\rho_i(u)\bigr)
\]
\struck{\uncertain{d'où}} \note{en interligne : « \uncertain{d'où} »}
\[
  \text{(31)}\qquad \sum_i f\bigl(\rho_i(vu)\bigr) \leq \sum_i f\bigl(\rho_i(v)\rho_i(u)\bigr)
\]
(\uncertain{la} \uncertain{somme}, \uncertain{finie} \uncertain{ou} \uncertain{infinie} \ill{} \uncertain{étant} \ill{} \uncertain{la} \uncertain{suite} $(i)$)

\note{Feuillet collé, couvrant le bas de la page :}
\ill{} $(\rho_i(vu))$ \uncertain{associée} \uncertain{à} \ill{} $u$, $v$, \uncertain{la} \ill{} \uncertain{majoration} \ill{} \uncertain{par} \uncertain{celle} \uncertain{qui} \ill{} \uncertain{si} \uncertain{on} \uncertain{avait} $\rho_i(vu) = \rho_i(u)\rho_i(v)$, \uncertain{donc} \uncertain{si} $u$ \uncertain{et} $v$ \uncertain{étaient} \uncertain{des} \uncertain{opérateurs} \uncertain{de} \uncertain{multiplication} \uncertain{dans} $\ell^2$ \uncertain{par} \uncertain{les} $(\rho_i(u))$ \uncertain{resp}. $(\rho_i(v))$. \struck{\ill{} \ill{} \ill{}} \note{en interligne, souligné : « \uncertain{Nous} \ill{} »} \struck{\ill{}} \ill{} \uncertain{d'où} \ill{} \uncertain{corollaire} \ill{} \uncertain{à} \ill{} \struck{\ill{}} \note{trois lignes biffées} \uncertain{corollaire} 5 \uncertain{du} \uncertain{Th}. 1, \uncertain{mais} \ill{} \uncertain{laissons} \ill{} \uncertain{au} \uncertain{lecteur}. \note{la phrase est reliée par un trait d'accolade}

\uncertain{Dans} \uncertain{l'inégalité} (30), \ill{} \uncertain{peut} \uncertain{faire} \uncertain{et} \ill{} \ill{} \uncertain{tend} \uncertain{vers} \ill{}, \uncertain{d'où} \ill{} \uncertain{par} (30), \ill{} \uncertain{et} \uncertain{on} \uncertain{obtient} \uncertain{à} \uncertain{la} \uncertain{limite}. \note{deux lignes barrées de traits ondulés}

\page{91}
\note{Feuillet étroit, pencilé 93 ; les formules (30) et (31) sont rappelées en tête, au crayon, au-dessus des numéros (27) et (28) qu'il cite ; sa moitié inférieure est un feuillet collé, portant le corollaire 4. Une longue note de côté occupe la marge gauche.}
\uncertain{En} \uncertain{effet}, 27 \note{« 30 » au-dessus, au crayon} \uncertain{est} \uncertain{un} \uncertain{cas} \uncertain{particulier} \uncertain{du} \uncertain{Th}. 2, \uncertain{et} 28 \note{« 31 » au-dessus, au crayon}, \uncertain{en} \uncertain{résulte} \uncertain{par} \uncertain{passage} \uncertain{à} \uncertain{la} \uncertain{limite} \uncertain{pour} $n \to +\infty$. \note{un astérisque au crayon rouge} \struck{\uncertain{En} \uncertain{particulier}} \struck{\uncertain{Prenons} \uncertain{en} \uncertain{particulier}, \uncertain{au} \uncertain{corollaire} 1, $f(r) = r^p$ \ill{} $p > 0$, \uncertain{on} \uncertain{obtient}} \note{deux lignes biffées} \uncertain{soient} $p, q, \struck{r}$ \uncertain{tels} \uncertain{que}
\[
  \text{(34)}\qquad \frac{1}{p} + \frac{1}{q} = \frac{1}{s} \qquad (p, q, s \geq 0)
\]
\note{le « (34) » est surchargé sur un autre numéro ; un $\sum$ biffé à gauche} \uncertain{posons} $f(\rho) = \rho^s$, \struck{\ill{} \uncertain{dans} \uncertain{le} \uncertain{corollaire} 1 \ill{}} \uncertain{l'inégalité} \struck{\ill{}} (31), \uncertain{donne}
\[
  \Bigl(\sum_{i=1}^\infty \rho_i(vu)^s\Bigr)^{1/s} \leq \Bigl(\sum_{i=1}^\infty \bigl(\rho_i(u)\rho_i(v)\bigr)^s\Bigr)^{1/s}
\]
\uncertain{Or} (\uncertain{inégalité} \uncertain{de} \uncertain{Hölder} \struck{\ill{}} \note{en interligne : « \uncertain{soit} \uncertain{bien} \uncertain{connue} »} \struck{\ill{} \uncertain{quels} \ill{} \uncertain{exposants} \ill{} \uncertain{conjugués}} \ill{} \uncertain{le} \uncertain{second} \uncertain{membre} \uncertain{est} \uncertain{majoré} \uncertain{par} \struck{$\bigl(\sum \rho_i(u)^s\bigr)$ \ill{}}
\[
  \Bigl(\sum \rho_i(u)^p\Bigr)^{1/p} \Bigl(\sum \rho_i(v)^q\Bigr)^{1/q} .
\]
\uncertain{On} \uncertain{trouve} \note{en interligne : « \ill{} »}

\emph{Corollaire 2}, (« \uncertain{Inégalité} \uncertain{de} \uncertain{Hölder} \struck{\uncertain{pour}} \uncertain{opérateurs} \ill{} »). \note{le titre entre guillemets est souligné, avec un trait rouge dessous ; un trait rouge dans la marge} \struck{\ill{} \ill{}} \uncertain{Soient} $E, F, G$ \uncertain{trois} \uncertain{espaces} \uncertain{de} \uncertain{Hilbert}, \uncertain{soit} $u$ \uncertain{une} \uncertain{op}. \uncertain{cpt} \uncertain{de} $E$ \uncertain{dans} $F$, $v$ \uncertain{une} \uncertain{op}. \uncertain{cpt} \uncertain{de} $F$ \uncertain{dans} $G$, \uncertain{alors} \uncertain{on} \uncertain{a}, \uncertain{pour} \struck{\uncertain{tout}} \ill{} \note{en interligne : « \ill{} entier »} $(p, q, s)$ \struck{\ill{}} \uncertain{satisfaisant} \uncertain{à} (34), \uncertain{et} \uncertain{tout} \uncertain{entier} $n > 0$ \note{l'énoncé est marqué d'un trait vertical dans la marge}
\[
  \text{(35)}\qquad \Bigl(\sum_{i=1}^n \rho_i(vu)^s\Bigr)^{1/s} \leq \Bigl(\sum_{i=1}^n \rho_i(u)^p\Bigr)^{1/p} \Bigl(\sum_{i=1}^n \rho_i(v)^q\Bigr)^{1/q}
\]
\uncertain{En} \uncertain{particulier}
\[
  \text{(36)}\qquad \Bigl(\sum \rho_i(vu)^s\Bigr)^{1/s} \leq \Bigl(\sum \rho_i(u)^p\Bigr)^{1/p} \Bigl(\sum \rho_i(v)^q\Bigr)^{1/q}
\]
\[
  \struck{N_s(vu) \leq N_p(u)\, N_q(v)} \qquad \|vu\|_s \leq \|u\|_p\, \|v\|_q \qquad \Bigl(\frac{1}{s} = \frac{1}{p} + \frac{1}{q}\Bigr)
\]
\note{les deux lignes de (36) sont encadrées ensemble au crayon rouge ; la forme en $N_s$ est biffée et récrite dessous avec les $\|\cdot\|$ ; une ligne biffée suit : « \struck{(\uncertain{où} \ill{} \uncertain{le} \uncertain{second} \uncertain{membre} \ill{} \uncertain{fini} \ill{})} »}

\uncertain{Le} \uncertain{cas} \uncertain{le} \uncertain{plus} \uncertain{important} \uncertain{de} \uncertain{l'inégalité} (37) \note{ainsi ; il s'agit de (36)} \uncertain{est} \uncertain{celui} \uncertain{où} $s = 1$, \uncertain{i.e.} \uncertain{où} $p$ \uncertain{et} $q$ \uncertain{sont} \uncertain{deux} « \uncertain{exposants} \uncertain{conjugués} » ; \struck{\uncertain{Nous} \uncertain{y} \uncertain{reviendrons} \ill{}} \struck{\ill{} \uncertain{inégalités} \uncertain{au} \uncertain{N°} \ill{}} \struck{\ill{} \uncertain{plus} \uncertain{frappante}} \note{trois lignes biffées}
\[
  \text{(38)}\qquad \|vu\|_1 \leq \|u\|_p\, \|v\|_{p'} \qquad \Bigl(\frac{1}{p} + \frac{1}{p'} = 1\Bigr)
\]
\note{formule encadrée au crayon rouge, le numéro cerclé} \struck{\ill{} \uncertain{l'inégalité} \uncertain{de} \uncertain{Hölder} \ill{}} (38) \uncertain{résulte} \uncertain{d'ailleurs} \uncertain{aussi} \uncertain{de} (32), \uncertain{grâce} \uncertain{à} \uncertain{l'inégalité} \uncertain{de} \uncertain{Hölder} \uncertain{ordinaire}. \note{la dernière ligne est barrée}

\marginal{(32) $\sum f\bigl(\rho_i(vu)\bigr) \leq \sum f\bigl(\rho_i(u)\rho_i(v)\bigr)$. \uncertain{Tenons} \uncertain{compte} \ill{} \ill{} \ill{} \uncertain{on} \uncertain{aura} $|\operatorname{Tr} vu| \leq \struck{\sum} \sum \rho_i(vu)$ (\uncertain{N°} 2, \uncertain{Th}. \ill{}) \ill{} \uncertain{formule} (21), \ill{} \uncertain{d'où} \ill{}. \uncertain{Prenons} \uncertain{en} \uncertain{particulier} \uncertain{$f(s) = s$}, \uncertain{on} \uncertain{obtient} \uncertain{de} \struck{(21)} \uncertain{N°} 1, \uncertain{Th}. 3, \ill{} $p = 1$, \struck{\ill{}} \ill{} $vu$ \uncertain{est} \uncertain{un} \uncertain{op}. \uncertain{de} \uncertain{Fredholm}. \uncertain{Par} \ill{} \struck{\uncertain{d'où} $|\operatorname{Tr}(vu)| \leq \sum |\lambda_i(vu)| \leq$} $|\operatorname{Tr} vu| \leq \sum \rho_i(u)\rho_i(v)$. \uncertain{Pour} \ill{} \emph{Corollaire 3} \uncertain{Soient} $E, F$ \uncertain{deux} \uncertain{espaces} \uncertain{de} \uncertain{Hilbert}, $u$ \uncertain{une} \uncertain{application} \uncertain{linéaire} \ill{} \uncertain{de} $E$ \uncertain{dans} $F$, \ill{} \uncertain{de} $F$ \uncertain{dans} $G$, \ill{} $\sum \rho_i(u)\rho_i(v)$ \ill{} \uncertain{on} \uncertain{a} \ill{} \uncertain{et} \ill{} (33) $|\operatorname{Tr} vu| \leq \sum \rho_i(u)\rho_i(v)$ \ill{} \uncertain{Ce} \uncertain{résultat} \ill{} \ill{} \uncertain{de} \uncertain{Fredholm} \ill{} (39).} \note{longue note de côté dans la marge gauche, courant sur toute la hauteur du feuillet et se poursuivant sur le feuillet collé ; les formules (32) et (33) y sont encadrées au crayon rouge, (32) cerclée ; un bloc biffé en zigzag en sépare les deux parties. Elle ajoute, entre le corollaire 2 et le corollaire 4, un corollaire 3 : l'inégalité de trace}

\note{Feuillet collé, moitié inférieure de la page ; un point d'interrogation au crayon rouge dans sa marge, et au crayon : « \uncertain{Est-ce} \uncertain{utile} ? \uncertain{d'où} \ill{} »}
\emph{Corollaire 4} \uncertain{Soient} $E, F, G$ \uncertain{trois} \uncertain{espaces} \uncertain{de} \uncertain{Hilbert} \uncertain{de} \uncertain{dimension} \uncertain{finie} $n$, $u$ \uncertain{une} \uncertain{op}. \ill{} \uncertain{de} $E$ \uncertain{dans} $F$, $v$ \uncertain{une} \uncertain{op}. \uncertain{de} $F$ \uncertain{dans} $G$, $f(\rho_1, \ldots, \rho_n)$ \uncertain{fonction} \struck{\ill{}} \uncertain{de} $n$ \uncertain{arguments} $\rho_i > 0$, \uncertain{telle} \uncertain{que} $f(e^{t_1}, \ldots, e^{t_n})$ \uncertain{soit} \uncertain{fonction} \uncertain{convexe} \uncertain{de} $(t_i) \in \mathbf{R}^n$. \uncertain{Sous} \uncertain{ces} \uncertain{conditions}, \note{en interligne : « \uncertain{si} $f$ \uncertain{est} \uncertain{définie} \uncertain{et} \uncertain{continue} \uncertain{pour} $\rho_i \geq 0$, \uncertain{ou} \uncertain{si} $u$ \uncertain{et} $v$ \uncertain{sont} \uncertain{inversibles} »} \uncertain{on} \uncertain{a} \note{l'énoncé est marqué d'un trait vertical dans la marge}
\[
  \text{(39)}\qquad f\bigl(\rho_1(vu), \ldots, \rho_n(vu)\bigr) \leq f\bigl(\rho_1(u)\rho_1(v), \ldots, \rho_n(u)\rho_n(v)\bigr)
\]
\uncertain{En} \uncertain{effet}, \struck{\ill{}} \uncertain{cela} \uncertain{résulte} \uncertain{du} \uncertain{Th}. 1, 2°, \uncertain{hyp}. \uncertain{b}), \ill{} \uncertain{en} \uncertain{vertu} \uncertain{de}
\[
  \text{(40)}\qquad \prod \rho_i(vu) = \struck{f(vu)}\; \prod \rho_i(u)\rho_i(v)
\]
\struck{\uncertain{En} \uncertain{effet}} \uncertain{Pour} \struck{\uncertain{établir}} \note{en interligne : « \uncertain{vérifier} »} (40), \uncertain{on} \struck{\uncertain{élève} \uncertain{au} \uncertain{carré}} \uncertain{peut} \uncertain{évidemment} \uncertain{supposer} \uncertain{que} $E = F = G$, \uncertain{alors}, \uncertain{en} \uncertain{vertu} \uncertain{de} \struck{\ill{}} (26), \uncertain{les} \uncertain{deux} \uncertain{membres} \uncertain{de} (40) \uncertain{sont} \uncertain{resp}. $|\det(vu)|$ \uncertain{et} $|\det u|\,|\det v|$ (\uncertain{formule} (26)), \uncertain{d'où} \uncertain{le} \uncertain{résultat}. \uncertain{Si} $u$ \uncertain{et} $v$ \uncertain{ne} \uncertain{sont} \uncertain{pas} \uncertain{inversibles}, \uncertain{mais} $f$ \uncertain{définie} \uncertain{et} \uncertain{continue} \uncertain{pour} $\rho_i \geq 0$, (39) \uncertain{se} \uncertain{ramène} \uncertain{par} \uncertain{passage} \uncertain{à} \uncertain{la} \uncertain{limite} \uncertain{à} \uncertain{partir} \uncertain{du} \uncertain{cas} \uncertain{précédent}.

\uncertain{Dans} \uncertain{des} \ill{} \uncertain{sur} \uncertain{la} \uncertain{comparaison} \uncertain{des} \uncertain{opérateurs} \ill{} \uncertain{antilinéaire} \note{deux lignes au crayon, très pâles, au bas du feuillet collé}

\page{92}
\emph{Th 4} (\emph{Troisième Th. de convexité dans les espaces de Hilbert}) \note{« Th 4 » souligné, le titre souligné ; l'énoncé est marqué d'un trait vertical dans la marge}

\uncertain{Soient} $E, F$ \uncertain{deux} \uncertain{espaces} \uncertain{de} \uncertain{Hilbert}, $u, v$ \uncertain{deux} \uncertain{opérateurs} \uncertain{compacts} \uncertain{de} $E$ \uncertain{dans} $F$, \uncertain{soit} \uncertain{enfin} $\varphi(\rho_1, \ldots, \rho_n)$ \uncertain{une} \uncertain{fonction} \struck{\uncertain{convexe} \uncertain{symétrique} \uncertain{croissante} \ill{}} \struck{\uncertain{continue} \uncertain{symétrique} \ill{}} \note{deux lignes biffées, avec en interligne « convexe symétrique croissante » récrit} \uncertain{de} $n$ \uncertain{arguments} $s_i \geq 0$, \struck{\uncertain{telle} \uncertain{que}} \note{en interligne : « Alors »} \struck{$f(e^{t_1}, \ldots, e^{t_n})$ \uncertain{soit} \uncertain{fonction} \uncertain{convexe} \uncertain{de} $(t_i) \in \mathbf{R}^n$. \uncertain{Alors}} \note{ligne biffée : l'hypothèse de convexité logarithmique est retirée}
\[
  \text{(41)}\qquad \varphi\bigl(\rho_1(u+v), \ldots, \rho_n(u+v)\bigr) \leq \varphi\bigl(\rho_1(u)+\rho_1(v), \ldots, \rho_n(u)+\rho_n(v)\bigr)
\]
\uncertain{En} \uncertain{effet}, \uncertain{c'est} \uncertain{un} \uncertain{cas} \uncertain{particulier} \uncertain{du} \uncertain{Th}. 1, 1°, \uncertain{hyp}. \uncertain{a}, \uncertain{compte} \uncertain{tenu} \uncertain{des} \uncertain{formules} \uncertain{du} \uncertain{N°} 1.

\struck{\emph{Corollaire.} \uncertain{Soient} $u$ \uncertain{et} $v$ \ill{} \uncertain{comme} \uncertain{ci-dessus}, \uncertain{soit} $\varphi(s)$ \uncertain{une} \uncertain{fonction} \uncertain{convexe} \uncertain{de} $s \geq 0$, \uncertain{alors} \uncertain{on} \uncertain{a} \uncertain{pour} \uncertain{tout} $n > 0$} \note{en interligne : « croissante »}
\[
  \struck{\text{(42)}\qquad \sum_1^n \varphi\bigl(\rho_i(u+v)\bigr) \leq \sum_1^n \varphi\bigl(\rho_i(u) + \rho_i(v)\bigr)}
\]
\note{le corollaire et la formule (42) sont entourés d'un cadre et barrés de grands traits}

\uncertain{Nous} \uncertain{serons} \uncertain{surtout} \uncertain{intéressés} \uncertain{par} \uncertain{le}

\emph{Corollaire 1} \uncertain{Soient} \struck{$u$ \uncertain{et} $v$ \uncertain{comme} \uncertain{ci-dessus}} \note{en interligne : « $E$ et $F$ deux espaces de Hilbert »}, \uncertain{et} \uncertain{soit} $N$ \emph{\uncertain{une} \uncertain{norme} \uncertain{de} \uncertain{Schatten}} \uncertain{sur} $\mathbf{R}^n$, \struck{\uncertain{alors} \uncertain{les} \uncertain{quantités}} \uncertain{posons}, \uncertain{pour} \struck{\uncertain{tout}} \uncertain{tout} \uncertain{opérateur} \uncertain{compact} $u$ \uncertain{de} $E$ \uncertain{dans} $F$
\[
  \text{(43)}\qquad \|u\|_N = N\bigl(\rho_1(u), \ldots, \rho_n(u)\bigr)
\]
\uncertain{Alors} \uncertain{on} \uncertain{a}
\[
  \text{(44)}\qquad \struck{\|\lambda u\|_N = |\lambda|\, \|u\|_N} \qquad \|u + v\|_N \leq \|u\|_N + \|v\|_N
\]
\note{la seconde ligne est encadrée au crayon rouge ; la première, biffée, est en partie prise dans le cadre} \struck{\ill{} \uncertain{est} \uncertain{donc} \uncertain{une} \uncertain{norme} \uncertain{sur} $(L_0(E,F))$ \uncertain{sous-espace} \ill{} \ill{} $\mathbf{R}^{(\mathbf{N})}$} \note{deux lignes biffées de traits en boucle}

\uncertain{En} \uncertain{effet}, $N(s_1, \ldots, s_n)$ \uncertain{pour} $s_i \geq 0$ \uncertain{satisfait} \uncertain{aux} \uncertain{conditions} \uncertain{de} (41), \uncertain{on} \uncertain{a} \uncertain{donc}
\[
  N\bigl(\rho_1(u+v), \ldots, \rho_n(u+v)\bigr) \leq N\bigl(\rho_1(u)+\rho_1(v), \ldots, \rho_n(u)+\rho_n(v)\bigr)
\]
\uncertain{et} \uncertain{le} \uncertain{second} \uncertain{membre} \uncertain{est} \uncertain{majoré} \uncertain{par} $N\bigl(\rho_1(u), \ldots, \rho_n(u)\bigr) + N\bigl(\rho_1(v), \ldots, \rho_n(v)\bigr)$ \uncertain{puisque} $N$ \uncertain{est} \uncertain{une} \uncertain{norme}, \uncertain{d'où} (44). \uncertain{En} \uncertain{particulier}

\emph{Corollaire 2} \struck{\uncertain{Soient}} \uncertain{si} $u, v \in \uncertain{L_0^{p}}(E,F)$ \note{l'exposant est peu net ; une flèche le désigne}, \struck{\uncertain{et} \ill{}} ($\struck{\chi}$ $1 \leq p \leq +\infty$). \struck{\ill{}} \struck{\uncertain{pour} \uncertain{tout} $n > 0$} \struck{$\|u+v\|_p \leq$}
\[
  \text{(45)}\qquad \Bigl(\sum_1^n \bigl(\rho_i(u+v)\bigr)^p\Bigr)^{1/p} \leq \Bigl(\sum_1^n \rho_i(u)^p\Bigr)^{1/p} + \Bigl(\sum_1^n \rho_i(v)^p\Bigr)^{1/p}
\]
\uncertain{d'où} \uncertain{par} \uncertain{passage} \uncertain{à} \uncertain{la} \uncertain{limite} \struck{\ill{}}
\[
  \text{(46)}\qquad \|u + v\|_p \leq \|u\|_p + \|v\|_p
\]
\note{formule encadrée au crayon rouge ; le reste de la page est blanc}

\section{Classes remarquables d'opérateurs dans l'espace de Hilbert}

\note{Le titre est de lui, en tête de la page 93 : « 5. Classes remarquables d'opérateurs dans l'espace de Hilbert ». La numérotation des formules repart de (1).}

\page{93}
5. Classes remarquables d'opérateurs dans l'espace de Hilbert

\uncertain{Nous} \uncertain{commençons} \uncertain{par} \uncertain{donner} \uncertain{un} \uncertain{exposé} \uncertain{de} \uncertain{la} \uncertain{théorie} \uncertain{de} \uncertain{Schatten}, \struck{\uncertain{qui} \uncertain{est}} \ill{} \uncertain{un} \uncertain{bon} \ill{} \uncertain{simplifiée} \ill{} \uncertain{par} \uncertain{la} \struck{\uncertain{lemme}} \note{en interligne : « inégalités de convexité »} \uncertain{de} \uncertain{Weyl}, \uncertain{et} \uncertain{en} \uncertain{y} \uncertain{apportant} \uncertain{quelques} \uncertain{compléments}. \note{un point d'interrogation au crayon rouge dans la marge, et deux traits rouges sous « exposé » et « apportant »} \struck{\ill{}} \uncertain{Rappelons} \uncertain{le}

\emph{Définition 1} \struck{\uncertain{Soit}} \uncertain{Une} \uncertain{norme} $N$ \uncertain{sur} $\mathbf{R}^n$ \uncertain{est} \uncertain{dite} \emph{\uncertain{norme} \uncertain{de} \uncertain{Schatten}} \uncertain{si} \uncertain{elle} \uncertain{est} \uncertain{symétrique}, \uncertain{satisfait} \uncertain{à} \note{le corps de la définition est marqué d'un trait vertical dans la marge}
\[
  \text{(1)}\qquad N(x_1, \ldots, x_n) = N(|x_1|, \ldots, |x_n|)
\]
\struck{\uncertain{et} \uncertain{si} \uncertain{elle}} 2. \note{« 2. » ou « i. »} \uncertain{elle} \uncertain{est} \uncertain{croissante} \struck{\ill{}} \uncertain{sur} \uncertain{l'ens}. \uncertain{des} $x_i \geq 0$, \uncertain{et} \struck{\uncertain{elle} \uncertain{satisfait}} \uncertain{est} \uncertain{normalisée} \uncertain{par} $N(1, 0, \ldots, 0) = 1$

\uncertain{Une} \uncertain{norme} $N$ \uncertain{sur} $\mathbf{R}^{(\mathbf{N})}$ \struck{$\mathbf{N}$ \ill{}} (\uncertain{suites} \ill{} \note{en interligne : « \ill{} »} \uncertain{infinies} \struck{\ill{}} \ill{} \uncertain{de} \uncertain{nombres} \uncertain{réels} \uncertain{n'ayant} \uncertain{qu'un} \uncertain{nombre} \uncertain{fini} \uncertain{de} \uncertain{termes} \ill{}) \uncertain{est} \uncertain{appelée} \uncertain{une} \uncertain{norme} \uncertain{de} \uncertain{Schatten}, \struck{\uncertain{pour} \uncertain{tout} \ill{} $n > 0$,} \uncertain{si} \uncertain{elle} \uncertain{induit} \uncertain{sur} \struck{\ill{}} \uncertain{sur} $\mathbf{R}^n$ \uncertain{une} \uncertain{norme} \uncertain{de} \uncertain{Schatten} ($\mathbf{R}^n$ \uncertain{est} \uncertain{ici} \uncertain{identifié} \struck{\uncertain{à} \uncertain{l'ensemble} \ill{}} \note{en interligne : « \uncertain{aux} \uncertain{suites} \ill{} »} \uncertain{de} $\mathbf{R}^{(\mathbf{N})}$ \ill{} \uncertain{nulles}, \ill{} \uncertain{à} \uncertain{partir} \uncertain{du} \uncertain{rang} $n+1$). \uncertain{Notons} \uncertain{qu'une} \uncertain{norme} \uncertain{de} \uncertain{Schatten} \uncertain{sur} $\mathbf{R}^n$ \uncertain{induit} \uncertain{sur} $\mathbf{R}^m$ ($m \leq n$) \uncertain{une} \uncertain{norme} \uncertain{de} \uncertain{Schatten}. \struck{\ill{} \ill{}} \note{en interligne : « \uncertain{Exemples} \uncertain{de} \uncertain{normes} \uncertain{de} »} \uncertain{Schatten} : \struck{$N_p((x_i)) = \bigl(\sum |x_i|^p\bigr)^{1/p}$ \ill{} $1 \leq p$} \note{formule biffée} \uncertain{les} \uncertain{normes} $N_p((x_i)) = \|(x_i)\|_p$, \uncertain{définies} \uncertain{pour} \struck{\ill{} $p = \infty$ \uncertain{par}} \note{ligne biffée}
\[
  \text{(2)}\qquad \|(x_i)\|_p = \Bigl(\sum |x_i|^p\Bigr)^{1/p} \qquad (1 \leq p < +\infty)
\]
\[
  \text{(3)}\qquad \|(x_i)\|_\infty = \operatorname{Sup}_i |x_i|
\]
\uncertain{Si} $N_0$ \uncertain{est} \uncertain{une} \uncertain{norme} \uncertain{de} \uncertain{Schatten} \struck{\ill{}} \uncertain{sur} $\mathbf{R}^n$, \uncertain{il} \uncertain{existe} \ill{} \uncertain{une} \uncertain{norme} \uncertain{de} \uncertain{Schatten} $N$ \uncertain{sur} $\mathbf{R}^{(\mathbf{N})}$ \note{en interligne : « qui »} \uncertain{l'induise} \struck{$N_0$ \ill{}}. \uncertain{Soit}, \uncertain{pour} \uncertain{tout} $x = (x_i) \in \uncertain{c_0}$ \note{« $\underline{C}_0$ », souligné ; incertain}, \ill{} \struck{\ill{}} \uncertain{soit} $(\uncertain{R_i}(x))$ \uncertain{la} \uncertain{suite} \uncertain{des} $|x_i|$ \struck{\ill{}} \note{en interligne : « \ill{} »} \uncertain{rangée} \uncertain{par} \uncertain{ordre} \uncertain{décroissant}, \uncertain{pour} \ill{}, \uncertain{posons} $N(x) = N_0\bigl(\rho_1(x), \ldots, \rho_n(x)\bigr)$. \uncertain{Il} \uncertain{est} \uncertain{trivial} \note{deux traits rouges sous les dernières lignes ; deux points d'interrogation au crayon rouge dans la marge}

\page{94}
\uncertain{que} $N(x)$ \uncertain{induit} \uncertain{sur} $\mathbf{R}^n$ \uncertain{la} \uncertain{fonction} $N_0(x)$, \uncertain{et} \uncertain{qu'on} \uncertain{a} $N(\lambda x) = |\lambda|\, N(x)$ \uncertain{et} $N(x) = 0 \Rightarrow x = 0$, \uncertain{il} \uncertain{faut} \uncertain{seulement} \uncertain{prouver} $N(x+y) \leq N(x) + N(y)$ (\uncertain{seule} \uncertain{chose} $=$ \uncertain{pour} $x, y \in c_0$) \note{« seule » souligné au crayon rouge ; le bloc est marqué d'un trait vertical dans la marge}. \uncertain{Pour} \uncertain{ceci}, \uncertain{on} \struck{\uncertain{regarde} $(x_i)$ \uncertain{et} $(y_i)$ \uncertain{comme}} \uncertain{remarque} \uncertain{que} \uncertain{si} $u_x$ \uncertain{est} \uncertain{l'opérateur} \uncertain{de} \uncertain{multiplication} \uncertain{par} $x$ \uncertain{dans} $\ell^2$ (\uncertain{qui} \uncertain{est} \uncertain{un} \uncertain{opérateur} \ill{} \uncertain{pour} $x \in c_0$), \uncertain{on} \uncertain{a} \struck{$u_{x+y} = u_x$} $\rho_i(u_x) = \rho_i(x)$, \uncertain{donc} \ill{} \uncertain{pour} \struck{\ill{}} $u_{x+y} = u_x + u_y$, \uncertain{et} \uncertain{l'inégalité} \uncertain{à} \uncertain{prouver} \uncertain{résulte} \uncertain{alors} \uncertain{du} \uncertain{N°} 4, \uncertain{Th}. 4, \uncertain{corollaire}. --- \uncertain{En} \uncertain{résumé}, \uncertain{les} \uncertain{normes} \struck{\ill{}} \note{en interligne : « de Schatten »} \struck{\uncertain{ainsi} \uncertain{associées} \uncertain{à} \uncertain{des} \uncertain{normes}} \note{ligne biffée} \uncertain{de} \uncertain{Schatten} $N_0$ \uncertain{sur} $\mathbf{R}^n$ \ill{} \uncertain{qu'il} \uncertain{existe} \uncertain{une} \uncertain{infinité} \note{en interligne : « \uncertain{continue} »} \uncertain{de} \uncertain{normes} \uncertain{de} \uncertain{Schatten} \uncertain{sur} $\mathbf{R}^{(\mathbf{N})}$ \uncertain{équivalentes} \uncertain{à} \uncertain{la} \uncertain{norme} \uncertain{induite} \uncertain{par} $c_0$.

\uncertain{Soit} $N$ \struck{\ill{}} \uncertain{une} \uncertain{norme} \uncertain{de} \uncertain{Schatten} \uncertain{sur} \struck{\ill{}} $\mathbf{R}^{(\mathbf{N})}$. \struck{\uncertain{Soit} $\ell^N$ \uncertain{l'espace} \uncertain{des} \uncertain{suites} \uncertain{complexes} $x = (x_i)$ \uncertain{telles} \uncertain{que} \ill{} \uncertain{soit} $N(x)$ \ill{}} \struck{$N(|x_1|, \ldots, |x_n|, 0, 0, \ldots)$ \uncertain{sont} \uncertain{avec} $n$, \ill{} \uncertain{finie}} \note{quatre lignes biffées, avec un « \uncertain{si} $n$ \uncertain{varie} » en interligne} \uncertain{Soit} $|x|_n = (|x_1|, \ldots, |x_n|, 0, \ldots)$. \struck{\uncertain{Comme} \ill{} \uncertain{si} $n$ \uncertain{varie} \ill{}} \struck{\uncertain{alors} $|x+y|_n \leq |x|_n + |y|_n$} \note{deux lignes biffées} $|x|_n$ \uncertain{est} \uncertain{une} \uncertain{suite} \uncertain{croissante} \uncertain{d'éléments} \uncertain{positifs} \uncertain{de} $\mathbf{R}^{(\mathbf{N})}$, \uncertain{posons}
\[
  \text{(4)}\qquad N(x) = \lim_n N(|x|_n)
\]
\struck{\uncertain{La} \uncertain{valeur} \ill{} $\geq 0$ \ill{}} \struck{\ill{}} \note{deux lignes biffées ; en interligne, cerclé : « \uncertain{à} $N(x)$, \uncertain{qui} \uncertain{coïncide} \uncertain{avec} \struck{$N(x)$} \uncertain{quand} $x \in \mathbf{R}^{(\mathbf{N})}$ »}
\[
  \text{(5)}\qquad N(\lambda x) = |\lambda|\, N(x), \qquad N(x+y) \leq N(x) + N(y),
\]
\struck{\ill{} $N(x) < +\infty$} \uncertain{la} \uncertain{première} \uncertain{formule}, \struck{\ill{}} \note{en interligne : « \uncertain{il} \uncertain{faut} \uncertain{prouver} \uncertain{la} \uncertain{deuxième} »} \struck{\ill{} \uncertain{i.e.} \ill{}} \struck{\ill{} 1, \uncertain{valeur} \ill{} \uncertain{s'il} \ill{}} \struck{\ill{} \uncertain{fini} \ill{}} \note{trois lignes biffées ; un mot souligné au crayon rouge} \uncertain{La} \uncertain{deuxième} \uncertain{formule} \uncertain{résulte} \uncertain{de} \uncertain{ce} \uncertain{que} \struck{$N(\ldots)$} $|x+y|_n \leq |x|_n + |y|_n$, \uncertain{d'où} $N(|x+y|_n) \leq N(|x|_n + |y|_n) \leq N(|x|_n) + N(|y|_n)$. \struck{\uncertain{Soit}} \uncertain{Soit} $\ell^N$ \uncertain{l'espace} \uncertain{des} \uncertain{suites} \uncertain{complexes} $x$ \uncertain{telles} \uncertain{que} $N(x) < +\infty$ \note{plusieurs mots biffés et récrits en interligne autour de cette phrase : « \uncertain{muni} \uncertain{de} », « \uncertain{c'est} \uncertain{un} \uncertain{espace} »} \ill{} \uncertain{pour} \ill{} $N(x)$ \struck{\uncertain{est}} \uncertain{une} \uncertain{norme} \uncertain{sur} $\ell^N$ \struck{\ill{}}. \struck{\uncertain{D'ailleurs} \ill{} $\mathbf{R}^{(\mathbf{N})} \subset \ell^N$ \ill{} $N(x)$ \ill{} \ill{}} \note{les trois dernières lignes de la page sont biffées de grands traits en boucle}

\page{95}
\struck{\ill{}, \uncertain{si} \ill{} \uncertain{ou} $N \leq$ \ill{}} \note{ligne biffée en tête de page}
\struck{(6) $\ell^1 \subset \ell^N \subset \ell^\infty$}
\struck{\uncertain{et} \uncertain{la} \uncertain{norme} \uncertain{de} $\ell^1$ \uncertain{est} \uncertain{plus} \uncertain{fine} \uncertain{que} \uncertain{celle} \uncertain{induite} \uncertain{par} $\ell^N$, \uncertain{qui} \uncertain{est} \uncertain{à} \uncertain{son} \uncertain{tour} \uncertain{plus} \uncertain{fine} \uncertain{que} \uncertain{celle} \uncertain{induite} \uncertain{par} \ill{}} \note{la formule et les trois lignes sont entourées d'un cadre et barrées d'un grand trait ondulé ; elles sont reprises plus bas comme (6)}

\uncertain{Si} $N'$ \uncertain{et} $N''$ \uncertain{sont} \uncertain{deux} \uncertain{normes} \uncertain{de} \uncertain{Schatten} \uncertain{telles} \uncertain{que} $N' \leq N''$, \uncertain{il} \uncertain{est} \uncertain{immédiat} \uncertain{que} $\ell^{N''} \subset \ell^{N'}$, \uncertain{et} \uncertain{la} \uncertain{norme} \uncertain{de} $\ell^{N'}$ \struck{$\ell^{N'} \subset \ell^{N''}$} \uncertain{est} \uncertain{moins} \uncertain{fine} \uncertain{que} \uncertain{celle} \uncertain{induite} \uncertain{par} $\ell^{N''}$. \struck{\uncertain{En} \uncertain{particulier},} \uncertain{En} \uncertain{particulier}, \uncertain{si} \uncertain{on} \uncertain{remarque} \uncertain{que} \uncertain{toute} \uncertain{norme} \uncertain{de} \uncertain{Schatten} \uncertain{est} \uncertain{comprise} \uncertain{entre} \uncertain{les} \uncertain{normes} \uncertain{de} \uncertain{Schatten} $N_1(x) = \sum_i |x_i|$ \uncertain{et} $N_\infty(x) = \operatorname{Sup}_i |x_i|$ \uncertain{et} \uncertain{que} \uncertain{les} \uncertain{espaces} $\ell^N$ \uncertain{associés} \uncertain{à} \uncertain{ces} \uncertain{deux} \uncertain{dernières} \uncertain{normes} \uncertain{sont} \struck{$\ell^{N_1}$ \uncertain{et} $\ell^{N_\infty}$} $\ell^1$ \uncertain{et} $\ell^\infty$, \ill{} \uncertain{on} \uncertain{voit} \uncertain{que}
\[
  \text{(6)}\qquad \ell^1 \subset \ell^N \subset \ell^\infty
\]
\uncertain{les} \uncertain{applications} \uncertain{identiques} \struck{\ill{}} $\ell^1 \to \ell^N$ \uncertain{et} $\ell^N \to \ell^\infty$ \uncertain{étant} \uncertain{de} \uncertain{norme} $\leq 1$. \uncertain{Notons} \ill{} \uncertain{l'inégalité} \note{la fin de la ligne est prise dans une longue boucle}
\[
  \text{(7)}\qquad \struck{\sum}\; N\Bigl(\sum_k x^k\Bigr) \leq \sum_k N(x^k)
\]
\uncertain{valable} \uncertain{pour} \uncertain{toute} \uncertain{suite} \struck{\ill{}} $(x^k)$ \uncertain{d'éléments} \ill{} \note{en interligne : « $x^k \in \ell^N$ »} \ill{} \uncertain{pour} \uncertain{laquelle} \ill{} \uncertain{convergente} \uncertain{dès} \uncertain{que} \uncertain{le} \uncertain{second} \uncertain{membre} \uncertain{est} \uncertain{fini}, \uncertain{i.e.} $\sum_k |x_i^k| < +\infty$. \uncertain{Il} \uncertain{suffit} \uncertain{en} \uncertain{effet} \uncertain{de} \uncertain{prouver} (7) \uncertain{si} \uncertain{le} \uncertain{second} \uncertain{membre} \uncertain{est} \uncertain{fini} ; \uncertain{alors} \uncertain{en} \uncertain{posant} $x = \sum_k x^k$, \uncertain{on} \uncertain{aura} $|x|_m \leq \sum_k |x^k|_m$ ; \struck{$N$ \uncertain{étant} \uncertain{une} \uncertain{fonction} \uncertain{croissante} \uncertain{sur} $\mathbf{R}^m$} \note{ligne biffée d'un trait} \uncertain{on} \uncertain{aura} $N(|x|_m) \leq N\bigl(\sum_k |x^k|_m\bigr)$, \uncertain{et} $N\bigl(\sum_k |x^k|_m\bigr) = \lim_{n \to \infty} N\bigl(\sum_{k=1}^n |x^k|_m\bigr)$, \uncertain{et} \uncertain{comme} $N\bigl(\sum_{k=1}^n |x^k|_m\bigr) \leq \sum_{k=1}^n N(|x^k|_m) \leq \sum_{k=1}^n N(x^k)$, \uncertain{on} \uncertain{aura} $N\bigl(\sum_k |x^k|_m\bigr) \leq \sum_k N(x^k)$, \uncertain{donc} \uncertain{pour} \uncertain{tout} $m$ $N(|x|_m) \leq \sum_k N(x^k)$, \uncertain{d'où} \uncertain{par} \uncertain{définition} \uncertain{de} $N(x)$ \uncertain{l'inégalité} (7). \note{la moitié inférieure de la page est traversée de grandes boucles et de traits obliques qui relient les lignes entre elles, sans les biffer : l'ordre de lecture est celui donné ici}

\page{96}
\uncertain{Le} \uncertain{membre} \uncertain{de} (7) \uncertain{que} $\ell^N$ \uncertain{est} \emph{\uncertain{complet}}. \uncertain{Il} \uncertain{suffit} \uncertain{de} \uncertain{prouver} \uncertain{qu'une} \uncertain{suite} $(x^k)$ \struck{\ill{} $\ell^N$} \uncertain{d'éléments} $x^k \in \ell^N$, \uncertain{telle} \uncertain{que} $\sum \|x^k\|_N < +\infty$, \uncertain{est} \uncertain{sommable} \uncertain{dans} $\ell^N$. \ill{} \uncertain{elle} \ill{} : \uncertain{à} \uncertain{fortiori} $\sum \|x^k\|_\infty < +\infty$, \uncertain{donc} \uncertain{la} \uncertain{série} \uncertain{est} \ill{} \uncertain{sommable} \uncertain{dans} $\ell^\infty$, \note{en interligne : « \uncertain{vers} \uncertain{un} $x \in \ell^\infty$ ; \uncertain{d'autre} \uncertain{part} »} \ill{} \uncertain{avec} \uncertain{soit} \ill{} (7), \ill{} $x \in \ell^N$. \uncertain{D'autre} \uncertain{part},
\[
  N\Bigl(x - \sum_{k=1}^n x^k\Bigr) = \struck{\ill{}}\; N\Bigl(\struck{x} \sum_{n+1}^\infty x^k\Bigr) \leq \sum_{n+1}^\infty N(x^k),
\]
\uncertain{qui} \uncertain{tend} \uncertain{vers} \uncertain{zéro} \struck{\ill{} \uncertain{pour} \ill{} $+\infty$} \note{fin de ligne biffée}, \uncertain{puisque} $\sum N(x^k)$ \uncertain{est} \uncertain{une} \uncertain{série} \uncertain{convergente}. \uncertain{Il} \uncertain{s'ensuit} \uncertain{que} \ill{} \uncertain{aussi} $\sum x^k = x$ \emph{\uncertain{dans} $\ell^N$}.
\marginal{\uncertain{Il} \uncertain{est} \uncertain{immédiat} \uncertain{de} \uncertain{vérifier} \uncertain{que} \ill{} \ill{} \uncertain{l'inégalité} \ill{} \ill{} $\ell^N$ \uncertain{est} \ill{}} \note{de côté dans la marge gauche, en regard du haut de la page, en cinq lignes courtes ; largement illisible}

\uncertain{Soit} $\ell_N$ \uncertain{l'adhérence} \uncertain{de} $\underline{\mathbf{R}^{(\mathbf{N})}}$ \uncertain{dans} $\ell^N$, \uncertain{comme} $\ell^N$ \uncertain{est} \uncertain{complet}, $\ell_N$ \uncertain{s'identifie} \uncertain{au} \uncertain{complété} \uncertain{de} $\mathbf{C}^{(\mathbf{N})}$ \uncertain{pour} \uncertain{la} \uncertain{norme} $N$. \uncertain{Il} \uncertain{est} \uncertain{encore} \uncertain{trivial} \uncertain{que} \uncertain{si} $N' \leq N''$, \uncertain{alors} $\ell_{N'} \subset \ell_{N''}$, \uncertain{et} \uncertain{l'application} \uncertain{identique} $\ell_{N'} \to \ell_{N''}$ \uncertain{est} \uncertain{de} \uncertain{norme} $\leq 1$. \note{ainsi sur la page : le sens de l'inclusion est l'inverse de celui de $\ell^{N''} \subset \ell^{N'}$ à la page précédente ; les indices $N'$, $N''$ sont peu nets} \uncertain{En} \uncertain{particulier}, \uncertain{les} \uncertain{espaces} $\ell_N$ \uncertain{associés} \uncertain{à} $N_1$ \uncertain{et} $N_\infty$ \uncertain{étant} \uncertain{resp}. $\ell^1$ \uncertain{et} $c_0$, \uncertain{on} \uncertain{obtient}
\[
  \text{(8)}\qquad \ell^1 \subset \ell_N \subset c_0
\]
\emph{Proposition 1} \struck{\ill{} $\ell^N \subset c_0$ \uncertain{si} \uncertain{et} \uncertain{seulement} \uncertain{si} $N$ \uncertain{est} \uncertain{équivalente} \uncertain{à} \ill{} : \uncertain{Notons} \uncertain{les} $N_\infty$ \ill{} \uncertain{propriétés} \uncertain{suivantes}} \note{le titre est souligné ; deux lignes biffées le suivent, et la formule qui suit est prise dans un grand trait en boucle}
\[
  \struck{\text{(9)}\qquad \ell_N = \ell^N \cap c_0}
\]
\struck{\uncertain{Il} \uncertain{suffit} \uncertain{de} \uncertain{prouver} \uncertain{que} $x \in \ell^N \cap c_0$ \uncertain{implique} $x \in \ell_N$. \ill{} \uncertain{En} \uncertain{effet} \ill{}} \note{deux lignes biffées, la seconde soulignée d'une ligne pointillée}

\uncertain{Notons} \uncertain{que} $\mathbf{R}^{(\mathbf{N})}$ \note{le $\mathbf{R}$ est repassé à l'encre épaisse} \uncertain{est} \uncertain{en} \uncertain{dualité} \uncertain{avec} \uncertain{lui-même} \struck{$=$} \uncertain{par} \uncertain{l'accouplement} $\langle x, y \rangle = \sum x_i y_i$. \uncertain{Cela} \uncertain{permet} \uncertain{de} \uncertain{faire} \uncertain{correspondre} \uncertain{à} \uncertain{toute} \uncertain{norme} $N$ \uncertain{sur} $\mathbf{R}^{(\mathbf{N})}$, \ill{} \uncertain{une} \emph{\uncertain{norme} \uncertain{polaire}}, \uncertain{notée} $N^{\circ}$, \uncertain{définie} \uncertain{par}
\[
  \text{(9)}\qquad N^{\circ}(y) = \operatorname{Sup}_{N(x) \leq 1} |\langle x, y \rangle|
\]
\note{le $y$ de $N^{\circ}(y)$ et de $\langle x, y \rangle$ est surchargé, écrit sur un $x'$}
\uncertain{Si} $N$ \uncertain{est} \uncertain{une} \uncertain{norme} \uncertain{de} \uncertain{Schatten}, \uncertain{on} \uncertain{vérifie} \uncertain{trivialement} \uncertain{que} $N^{\circ}$ \uncertain{est} \uncertain{une} \uncertain{norme} \uncertain{de} \uncertain{Schatten}, \uncertain{et} \struck{\uncertain{si} $x \in \ell_N$, $x' \in \ell^{N^{\circ}}$, \uncertain{alors} \ill{} \ill{} $\ell^{N'}$, \ill{} $\sum x_i x_i'$ \uncertain{est} \ill{}, \ill{} \uncertain{et} \ill{}} \note{les deux dernières lignes sont biffées et reliées par une accolade au « et »}
\marginal{\uncertain{De} \uncertain{plus}, \uncertain{clairement} $= (N^{\circ})^{\circ} = N$, \ill{} \uncertain{vérifier} \ill{}} \note{de côté dans la marge gauche, au bas de la page, en regard de la norme polaire ; sa fin est biffée}

\page{97}
\[
  \text{(10)}\qquad \sum |x_i x_i'| \leq N(x)\, N'(x')
\]
\note{le « (10) » est surchargé ; $N'$ désigne ici la norme polaire $N^{\circ}$ de la page précédente, et l'écriture $N'$ est celle de toute la suite}
\uncertain{En} \uncertain{effet}, \uncertain{on} \uncertain{a} $\sum_1^n |x_i x_i'| = \langle |x|_n, |x'|_n \rangle \leq N(|x|_n)\, N'(|x'|_n) \leq N(x)\, N'(x')$, \uncertain{d'où} (10). \uncertain{Il} \uncertain{en} \uncertain{résulte} \uncertain{un} \uncertain{accouplement} \uncertain{naturel} \uncertain{entre} $\ell^N$ \uncertain{et} $\ell^{N'}$ \note{« naturel » souligné} \struck{\ill{} \ill{} \ill{}} \note{ligne biffée, avec en interligne « \uncertain{i.e.} \uncertain{une} \uncertain{application} \uncertain{linéaire} \uncertain{continue} », « \uncertain{de} \uncertain{norme} $\leq 1$ »} $\ell_N$ \uncertain{et} $\ell^{N'}$, \struck{\ill{} \uncertain{par} \uncertain{lequel} \uncertain{le} \uncertain{dual} \uncertain{de} $\ell_N$ \uncertain{s'identifie}} \note{ligne biffée, cerclée à sa fin ; deux traits rouges « !! » dans la marge en regard} \uncertain{est} \uncertain{identique} \uncertain{à} $\ell^{N'}$ \note{souligné} (\uncertain{avec} \uncertain{sa} \uncertain{norme}) \uncertain{pour} \ill{} \uncertain{accouplement}. \uncertain{En} \uncertain{effet}, \struck{\ill{}} \uncertain{il} \uncertain{suffit} \uncertain{de} \uncertain{montrer} \uncertain{qu'une} \uncertain{forme} \uncertain{linéaire} \uncertain{continue} $X'$ \uncertain{sur} $\ell_N$ \uncertain{est} \uncertain{définie} \uncertain{par} \uncertain{un} $x'$ \uncertain{élément} \uncertain{de} $\ell^{N'}$ \struck{\ill{}} \uncertain{tel} \uncertain{que} $N'(x') \leq \|X'\|$. \uncertain{En} \uncertain{effet}, $X'$ \struck{\uncertain{est} \ill{}} \note{en interligne : « aussi »} \uncertain{il} \uncertain{suffit} \uncertain{de} \uncertain{le} \uncertain{voir} \uncertain{pour} \uncertain{sa} \struck{\ill{}} \uncertain{restriction} \uncertain{à} $\ell^1$, \uncertain{alors} $X'$ \struck{\uncertain{est} \uncertain{définie} \uncertain{par} \uncertain{un}} \note{en interligne : « \uncertain{est} »} $x' \in \ell^\infty$, \struck{\ill{}} \ill{} $N'(|x'|_n) \leq \|X'\|$, \uncertain{donc} \ill{} $N'(x') \leq \|X'\|$, \uncertain{et} \uncertain{le} \uncertain{résultat}. --- \uncertain{Notons} \uncertain{qu'en} \uncertain{général}, \uncertain{le} \uncertain{dual} \uncertain{de} $\ell^N$ \uncertain{n'est} \uncertain{pas} $\ell^{N'}$, \uncertain{car} $\ell_N$ \uncertain{n'est} \uncertain{pas} \uncertain{dense} \uncertain{dans} $\ell^N$. \uncertain{Ex} : $N = N_\infty$. \uncertain{Notons} \uncertain{que}

\emph{Prop} \note{dans la marge, souligné, avec un point d'interrogation au crayon rouge en dessous} \uncertain{pour} \uncertain{que} $\ell_N$ \uncertain{soit} \uncertain{réflexif}, \uncertain{il} \uncertain{faut} \uncertain{et} \uncertain{il} \uncertain{suffit} \uncertain{que} \struck{$\ell^N$ \ill{} $\ell_{N'}$ \ill{} $\ell^{N'}$ \uncertain{le} \uncertain{soit}} \note{en interligne : « \uncertain{il} \uncertain{faut} \uncertain{et} \uncertain{il} \uncertain{suffit} \uncertain{qu'on} \uncertain{ait} »} \struck{\uncertain{En} \uncertain{effet},} \uncertain{on} \uncertain{ait} $\ell_N = \ell^N$ \uncertain{et} $\ell_{N'} = \ell^{N'}$ \note{cette ligne est soulignée ; l'énoncé est souligné en entier} \struck{$\ell^N$ \uncertain{réflexif} \uncertain{implique} $\ell_N$ \uncertain{réflexif} (\uncertain{sous-espace} \uncertain{fermé}) \uncertain{qui} \uncertain{implique} $\ell^{N'}$ \uncertain{réflexif} (\uncertain{dual} \uncertain{de} $\ell_N$) \uncertain{qui} \uncertain{implique} $\ell_{N'}$ \uncertain{réflexif} \uncertain{et} \uncertain{donc} \ill{} $\ell^N$ \uncertain{réflexif}, \ill{}} \note{quatre lignes biffées d'un trait chacune, avec des mots en interligne : « \uncertain{de} \uncertain{même} \uncertain{que} », « \uncertain{d'un} \uncertain{réflexif} »} \uncertain{qui} \uncertain{admise} \uncertain{la} \uncertain{démonstration} \uncertain{circulaire} \uncertain{de} \uncertain{l'équivalence} \uncertain{des} 4 \uncertain{conditions}. \uncertain{Alors} $\ell_{N'}$, \uncertain{qui} \uncertain{est} \uncertain{un} \uncertain{sous-espace} \note{en interligne : « \uncertain{il} \uncertain{faut} \uncertain{donc} \uncertain{fermé} »} \struck{\uncertain{du} \uncertain{dual}} $\ell^{N'}$ \uncertain{de} $\ell_N$, \uncertain{est} \uncertain{réflexif} ; \note{« réflexif » ou « isomorphe », souligné au crayon rouge, avec un point d'interrogation rouge dans la marge} $\ell^{N'}$, \uncertain{il} \uncertain{par} \ill{} $=$ \ill{} \uncertain{on} \uncertain{aura} $\ell_N = \ell^N$. \uncertain{Réciproquement}, \struck{\uncertain{s'il} \uncertain{en} \uncertain{est} \uncertain{ainsi}}, \uncertain{alors} \uncertain{le} \uncertain{dual} \uncertain{de} $\ell_N$ \uncertain{est} $\ell^{N'} = \ell_{N'}$, \uncertain{et} \uncertain{le} \uncertain{dual} \uncertain{de} \uncertain{ce} \uncertain{dernier} \uncertain{est} $\ell^N = \ell_N$, \uncertain{donc} $\ell_N$ \uncertain{coïncide} \uncertain{avec} \uncertain{son} \uncertain{bidual}, \uncertain{et} \uncertain{est} \uncertain{donc} \uncertain{réflexif}.

\uncertain{Soit} $N$ \uncertain{une} \uncertain{norme} \uncertain{de} \uncertain{Schatten} \uncertain{sur} $\mathbf{R}^{(\mathbf{N})}$ \struck{\uncertain{non} \uncertain{équivalente} \uncertain{à} $N_\infty$} \note{en interligne, biffé, avec au-dessus « \uncertain{non} \uncertain{équivalente} \uncertain{à} $N_\infty$ » récrit}. \uncertain{Soient} $E, F$ \uncertain{deux} \uncertain{espaces} \uncertain{de} \uncertain{Hilbert}, \struck{$N$} \uncertain{on} \uncertain{désigne} \uncertain{par} $L^N(E,F)$ \note{le $L$ est repassé, l'exposant $N$ surchargé d'un indice biffé} \struck{\uncertain{l'espace}} \uncertain{sous-espace} \uncertain{de} $L_0(E,F)$ \uncertain{formé} \uncertain{des} $u$ \uncertain{tels} \uncertain{que} $(\rho_i(u)) \in \ell^N$, \uncertain{i.e.} \uncertain{tels} \uncertain{que}
\[
  \text{(11)}\qquad N(u) = N\bigl((\rho_i(u))\bigr) < +\infty
\]

\page{98}
\struck{\uncertain{Si} $N$ \ill{} \ill{}} \struck{\uncertain{Sous} \uncertain{ces} \uncertain{conditions},} \note{deux lignes biffées en tête de page} \uncertain{Alors} \uncertain{la} \uncertain{fonction} $u \mapsto N(u)$ \uncertain{sur} $L^N(E,F)$ \uncertain{est} \uncertain{une} \emph{\uncertain{norme}}. \uncertain{En} \uncertain{effet}, \uncertain{on} \uncertain{a} $N(u+v) = \lim_{n \to \infty} N\bigl(|\rho_i(u+v)|_n\bigr)$ \struck{\ill{}}, \uncertain{le} \uncertain{second} \uncertain{membre} \uncertain{est} \uncertain{majoré} \uncertain{par} $N\bigl(|\rho_i(u)|_n\bigr) \struck{\ill{}} + N\bigl(|\rho_i(v)|_n\bigr) \struck{\ill{}}$ (\uncertain{N°} 4, \uncertain{formule} (45)), \uncertain{donc} \uncertain{par} $N(u) + N(v)$. \uncertain{On} \uncertain{voit} \uncertain{trivialement} $N(\lambda u) = |\lambda|\, N(u)$, \uncertain{et} $N(u) = 0 \Rightarrow u = 0$. \uncertain{L'espace} $L^N(E,F)$ \uncertain{est} \uncertain{donc} \uncertain{un} \uncertain{espace} \uncertain{normé} ; \struck{\uncertain{prouvons}} \note{en interligne : « soit »} \struck{\uncertain{qu'il} \uncertain{est} \uncertain{complet}} \note{« complet » souligné d'un trait double, la ligne biffée}. $L_N(E,F)$ \uncertain{le} \uncertain{sous-espace} \uncertain{normé} \uncertain{formé} \uncertain{des} $u$ \uncertain{tels} \uncertain{que} $(\rho_i(u)) \in \ell_N$. \uncertain{Je} \uncertain{dis} \uncertain{que} $L_N(E,F)$ \uncertain{est} \uncertain{l'adhérence} \uncertain{de} $\bar{E} \otimes F$ \uncertain{dans} $L^N(E,F)$. \uncertain{En} \uncertain{effet}, \uncertain{Montrons} \uncertain{que} $\sum \rho_i\, \bar{e}_i \otimes f_i$ \uncertain{est} \uncertain{limite} \uncertain{des} \uncertain{sommes} \uncertain{partielles} $\sum_1^n \rho_i\, \bar{e}_i \otimes f_i$, \uncertain{dans} \uncertain{le} \uncertain{cas} \uncertain{où} $N\bigl(\sum_{n+1}^\infty \rho_i\, \bar{e}_i \otimes f_i\bigr) = N(\rho_{n+1}, \rho_{n+2}, \ldots) \to 0$ \uncertain{pour} $n \to +\infty$, \uncertain{ce} \uncertain{qui} \uncertain{démontre} \uncertain{donc}, \uncertain{que} \uncertain{dans} $\ell_N$ \uncertain{les} \uncertain{opérations} $x \mapsto t_n(x) = (0, \ldots, 0, x_{n+1}, x_{n+2}, \ldots)$ \note{les $n$ zéros sont marqués d'une accolade « $n$ »} \uncertain{tendent} \uncertain{vers} $0$ \uncertain{simplement}. \uncertain{Les} \uncertain{deux} \uncertain{formes} \uncertain{sont} \uncertain{équivalentes}, \uncertain{où} \uncertain{tout} \uncertain{ceci} \uncertain{est} \uncertain{zéro} \uncertain{dans} \ill{}, \uncertain{on} \uncertain{peut} \uncertain{le} \uncertain{avoir}. \struck{\ill{} $\ell^N(E,F)$} \note{fin de ligne biffée}

\uncertain{Le} \uncertain{dual} \uncertain{de} $L_N(E,F)$ \uncertain{s'identifie} \uncertain{à} \uncertain{un} \uncertain{sous-espace} \uncertain{du} \uncertain{dual} \uncertain{de} $L_{N_1}(E,F) = L^1(E,F)$, \uncertain{donc} \uncertain{de} $L(F,E)$. \uncertain{Il} \uncertain{résulte} \uncertain{que} $N$ \uncertain{n'est} \uncertain{pas} \uncertain{équivalente} \uncertain{à} $N_1$, \uncertain{donc} $N'$ \uncertain{n'est} \uncertain{pas} \uncertain{équivalente} \uncertain{à} $N_\infty$. \uncertain{Je} \uncertain{dis} \uncertain{que}
\[
  \bigl(L_N(E,F)\bigr)' = L^{N'}(F,E) .
\]
\uncertain{En} \uncertain{premier} \uncertain{lieu}, \uncertain{si} $u \in L_N(E,F)$, $v \in L^{N'}(F,E)$, \uncertain{donc} $vu \in L^1(E,E)$, \uncertain{avec}
\[
  \|vu\|_1 \leq N(u)\, N'(v)
\]
\note{tout ce bloc, du « Le dual » à la formule, est entouré d'un cadre et traversé de trois longs traits obliques ; le bloc est repris à la page suivante}
\marginal{(12) $N(u^{*}) = N(u)$ ; (13) $N(vu) \leq \|v\|\, N(u)$, $N(uv) \leq \|v\|\, N(u)$ ; (14) $N(u) = \operatorname{Sup} \ill{}$ pour $u' \in \bar{F}' \otimes E$, $N^{\circ}(u') \leq 1$ \ill{} $L^N(E,F)$ \uncertain{si} $N = N^{\circ}$ \ill{} \uncertain{définition} \uncertain{de} $L^N(E,F)$ ; ($L^{N_\infty}(E,F) = L(E,F)$ \ill{} $L_{N_\infty}(E,F)$ \ill{}) \ill{}} \note{trois notes de côté dans la marge gauche : les formules (12), (13) en haut, la formule (14), cerclée et encadrée, au milieu, et une parenthèse au bas sur les cas $N = N_\infty$ ; l'argument de (14) est en partie perdu}

\uncertain{Soient} $N, N', N''$ \uncertain{des} \uncertain{normes} \uncertain{de} \uncertain{Schatten} \note{la ligne est soulignée d'un long trait} \uncertain{telles} \uncertain{que}
\[
  \text{(\uncertain{15})}\qquad N(xy) \leq N'(x)\, N''(y)
\]
\note{formule encadrée ; le numéro est surchargé, (13) ou (15)} \uncertain{pour} \struck{\uncertain{tout}} $x, y \in \mathbf{R}^{(\mathbf{N})}$. \uncertain{Alors} \uncertain{la} \uncertain{même} \uncertain{formule} \uncertain{est} \uncertain{valable} \uncertain{pour} $x \struck{y} \in \ell^{N'}$, $y \in \ell^{N''}$. \uncertain{Par} \uncertain{suite}, \uncertain{si} $u \in L^{N'}(E,F)$, $v \in L^{N''}(F,G)$, \uncertain{on} \uncertain{a} $vu \in L^N(E,G)$, \uncertain{avec}
\[
  \text{(16)}\qquad N(vu) \leq N'(u)\, N''(v)
\]
\note{formule encadrée, au bas de la page}

\page{99}
\uncertain{Ce} \uncertain{résultat} \uncertain{résulte} \note{en interligne au-dessus, cerclé : « \uncertain{d'où}, $L$ \ill{} \uncertain{si} $N' \neq N_\infty$, $N^{\circ} \neq N_\infty$, »} \struck{\uncertain{en} \uncertain{admission} \uncertain{sur} \uncertain{la} \uncertain{caractérisation} \ill{}} \note{en interligne : « \uncertain{de} \uncertain{la} \uncertain{définition} »} \uncertain{précédente}, \uncertain{et} \uncertain{des} \uncertain{formules} $\rho_{i}(uv) \struck{-} \ill{} \leq \bigl(\rho_i(u)\rho_j(v)\bigr) \ill{}$ \note{formule très serrée en fin de ligne : une majoration de $\rho_{i+j}(uv)$ par $\rho_i(u)\rho_j(v)$, incertaine} (\uncertain{N°} 2, \uncertain{formule} (37)). --- \uncertain{En} \uncertain{particulier}, \struck{\ill{} \uncertain{à} \uncertain{la} \uncertain{fois}} \note{en interligne : « \ill{} »} \struck{\ill{}} \uncertain{aux} \uncertain{normes} \uncertain{polaires} $N$ \uncertain{et} $N^{\circ}$, \uncertain{on} \uncertain{trouve}
\[
  \text{(17)}\qquad \|uv\|_1 \leq N(u)\, N^{\circ}(v)
\]
\note{formule encadrée ; le « (17) » est cerclé et barré d'une croix ; en fin de ligne, biffé : « \struck{\uncertain{en} \uncertain{tenant} \uncertain{compte}} »}
\struck{(\ill{}) $N(uv) \leq \|u\|\, N(v)$, $N(uv) \leq \|v\|\, N(u)$} \note{ligne encadrée et biffée d'un trait ondulé, avec un mot cerclé et biffé à sa gauche ; les formules sont celles de la note marginale (13) de la page précédente}

\uncertain{Le} \uncertain{dual} \uncertain{de} $L_N(E,F)$ \uncertain{s'identifie} \uncertain{à} \uncertain{un} \uncertain{sous-espace} \uncertain{du} \uncertain{dual} \uncertain{de} $L_{N_1}(E,F) = L^1(E,F)$, \uncertain{donc} \uncertain{de} $L(F,E)$. \note{un trait rouge dans la marge et sous « $L_{N_1}(E,F) = L^1(E,F)$ »} \uncertain{Ce} \uncertain{sous-espace} \uncertain{contient} $L^{N^{\circ}}(F,E)$, \uncertain{et} \uncertain{l'application} $L^{N^{\circ}}(F,E) \to \bigl(L_N(E,F)\bigr)'$ \uncertain{est} \uncertain{un} \uncertain{morphisme} \uncertain{métrique} \note{« métrique » souligné}. \uncertain{Montrons} \uncertain{que} \uncertain{c'est} \emph{\uncertain{sur}}. \uncertain{On} \uncertain{peut} \uncertain{supposer} \ill{} $N \neq N_1$, \struck{\uncertain{prouvons}} \note{en interligne : « \uncertain{Soit} »} \ill{} \note{en interligne : « \uncertain{admet} », « \uncertain{i.e.} \uncertain{a} »} $\Phi \in \bigl(L_N(E,F)\bigr)'$, \ill{} $v \in L(F,E)$, \uncertain{i.e.} \uncertain{a} $v \in L^{N^{\circ}}(F,E)$. \uncertain{On} \ill{} \ill{} \struck{\ill{} \uncertain{d'où} $F$ \ill{}, $E$. \ill{} $v^{*} v$ \ill{}} \struck{\ill{} $|\operatorname{Tr} vu| \leq M\, N(u)$, \uncertain{pour} $u \in E' \otimes F$, \ill{}} \struck{$v v^{*} \in L(E,F)$ \ill{} \ill{} \ill{}} \note{trois lignes biffées d'un trait chacune} \struck{\ill{}} $v = U|v|$, \uncertain{il} \uncertain{suffit} \uncertain{de} \uncertain{prouver} \uncertain{que} $|v| \in L^{N^{\circ}}(E)$ \note{un $F$ ou $E$ surchargé après la parenthèse}, \uncertain{or} $|v| = V v$, \uncertain{donc} \ill{} \uncertain{la} \uncertain{même} \uncertain{condition} : $v$ ($|\operatorname{Tr}(vu)| \leq M\, N(u)$ \uncertain{pour} $u \in E' \otimes F$, $M$ \uncertain{constante}). \uncertain{On} \uncertain{est} \uncertain{ramené} \uncertain{à} : $F = E$, $v \geq 0$. \ill{} \uncertain{et} \uncertain{quel} \ill{} \uncertain{il} \uncertain{existe} \uncertain{un} \uncertain{projecteur} $w$ \uncertain{tel} \uncertain{que} $\struck{\ill{}}\, w v = P$, \uncertain{projecteur} \uncertain{hermitien} \uncertain{de} \uncertain{rang} \uncertain{fini}, \struck{\ill{} \ill{}} \ill{} \uncertain{aussi} \ill{} $w \in \bigl(L_N(E,E)\bigr)'$, \ill{} \ill{} \uncertain{compatible}, \struck{\ill{} \ill{}} \note{en interligne : « \ill{} »} \struck{\ill{} \ill{} $\|u\|_1 \leq P$, \ill{}} \note{deux lignes biffées, avec « $\operatorname{Tr} P =$ » barré} \ill{} \uncertain{soit} \ill{} $P$, \ill{} \uncertain{soit} \ill{} \ill{} $1 \in (\ell_N)' = \ell^{N^{\circ}}$, \ill{} \ill{} \ill{}. \ill{} \ill{} $v = \sum \rho_i\, \bar{e}_i \otimes e_i$, \uncertain{et} \uncertain{on} \uncertain{trouve} \note{« trouve » souligné au crayon rouge} \ill{} \ill{} \uncertain{que} $(\rho_i) \in (\ell_N)' = \ell^{N^{\circ}}$, \uncertain{d'où} \uncertain{la} \uncertain{conclusion}. --- \uncertain{Cette} \uncertain{conséquence} \note{souligné au crayon rouge ; un point d'interrogation rouge dans la marge} \ill{} \ill{} $L^{N^{\circ}}(E,F)$ \ill{} \uncertain{F} \ill{} \note{souligné au crayon rouge} \ill{}, \ill{} \ill{} $L^N(E,F)$ \ill{} \ill{} \ill{} $L_N(E,F)$ \note{la dernière ligne, soulignée au crayon rouge, est en partie perdue}
\marginal{$N^{\circ}$ \ill{} $N' = $ \ill{} \ill{} $N^{\circ}$ \ill{} \ill{} $N'' \neq$ \ill{} \ill{} $N(\rho_i)$ \ill{} $N(u) = \lim$ \ill{} \ill{} $P$ \ill{}} \note{une note au crayon, de côté dans la marge gauche, en regard du haut de la page, en une dizaine de lignes courtes ; elle porte sur $N^{\circ}$ et $N'$ et reste illisible même redressée}

\page{100}
\emph{Corollaire 1} $u \in L^N(E,F)$ \uncertain{si} \uncertain{et} \uncertain{seulement} \uncertain{si} $|u| \in L^N(E,F)$, \uncertain{et} \note{l'énoncé est marqué d'un trait vertical dans la marge}
\[
  N(u) = N(|u|)
\]
$u \in L_N(E,F)$ \uncertain{si} \uncertain{et} \uncertain{seulement} \uncertain{si} $|u| \in L_N(E,F)$.

\struck{\emph{Corollaire 2} \ill{} \uncertain{un} \uncertain{opérateur} \uncertain{compact} \uncertain{sur} \uncertain{l'espace} \uncertain{de} \uncertain{Hilbert} \ill{} $u^{*} u$ \ill{} $E$. \uncertain{En} \uncertain{effet}, $(\rho_i(u))$ \ill{} \uncertain{fonction} \uncertain{croissante} \ill{}} \note{trois lignes barrées de grands zigzags}

\emph{Corollaire 2} \uncertain{Si} $u = h + ik$, $h$ \uncertain{et} $k$ \uncertain{hermitiens}, \uncertain{alors}
\[
  u \in L_N(E) \iff h, k \in L_N(E), \qquad u \in L^N(E) \iff h, k \in L^N(E) .
\]

\emph{Remarque} \uncertain{Soit} $p$ \uncertain{une} \uncertain{norme} \uncertain{sur} $E' \otimes F$ \uncertain{satisfaisant} \uncertain{aux} \uncertain{conditions} : a) $p(Vu) = p(u)$, $p(uU) = p(u)$ \uncertain{pour} $V$ \uncertain{unitaire} \uncertain{dans} $F$, $U$ \uncertain{unitaire} \uncertain{dans} $E$ \note{les deux égalités sont écrites l'une sous l'autre ; « $U$ » est biffé puis récrit, « $E$ » ajouté au-dessus de « $F$ »} ; b) $p(x' \otimes y) = p(x')\, p(y)$ \note{ainsi sur la page : $p$ désigne aussi les normes de $E'$ et de $F$}. \uncertain{Supposons} $E$ \uncertain{et} $F$ \uncertain{de} \uncertain{dimension} \uncertain{infinie} \uncertain{pour} \uncertain{fixer} \uncertain{les} \uncertain{idées}, \uncertain{soient} $(e_i)$, $(f_i)$ \uncertain{des} \uncertain{systèmes} \uncertain{orthonormaux} \uncertain{infinis} \uncertain{dans} $E$ \uncertain{et} $F$, \ill{} \uncertain{considérons} $\mathbf{R}^{(\mathbf{N})} \to E' \otimes F \subset L(E,F)$ \note{le $\mathbf{N}$ est souligné, avec un $\varphi$ ou $\lambda$ au-dessus} \uncertain{par} $(\lambda_i) \mapsto \sum \lambda_i\, \bar{e}_i \otimes f_i$. \uncertain{Alors} $p(\varphi(\lambda))$ \uncertain{sur} $\mathbf{R}^{(\mathbf{N})}$ \uncertain{est} \uncertain{une} \uncertain{norme}, \uncertain{induisant} \uncertain{sur} $\mathbf{R}^n$ \uncertain{des} \uncertain{normes} \struck{\uncertain{de} \uncertain{Schatten}} \uncertain{invariantes} \uncertain{par} $\mathfrak{S}_n$ (\uncertain{condition} a) \ill{} $p$) \uncertain{normalisées} (\uncertain{condition} b) \ill{} \uncertain{dépendant} \ill{} \uncertain{que} \uncertain{de} $(|\lambda_i|)$ (\uncertain{condition} a)). \struck{\uncertain{Elle} $N$ \uncertain{est} \uncertain{donc} \uncertain{une}} \note{souligné et biffé} \uncertain{norme} \uncertain{de} \uncertain{Schatten}. \uncertain{On} \uncertain{vérifie} \uncertain{alors}, \uncertain{considérant} (\uncertain{condition} a)) \uncertain{que} \uncertain{l'on} \uncertain{a} $N(u) = p(u)$ \uncertain{pour} \struck{\uncertain{tout}} $u \in E' \otimes F$. \uncertain{On} \uncertain{a} \uncertain{ainsi} \uncertain{obtenu} \uncertain{une} \uncertain{caractérisation} \uncertain{axiomatique} \uncertain{des} \uncertain{normes} $N(u)$ \uncertain{sur} \uncertain{les} \uncertain{espaces} \struck{$L(E,F)$} $E' \otimes F$. \note{le reste de la page est blanc}

\end{document}
